% Appendix C --- the complete proof of thm:wedge-axis (Q-0312, review flags 3 and 30). The review
% asked for complete proofs of the section's main theorem and of every statement its proof spends
% that the paper only sketched. This appendix transcribes the blueprint's proofs of record
% (blueprint/src/parts/05-*.tex), in the order of their dependence, with the paper's pointers
% (`Lemma~\ref{...}` for the blueprint's bare `\ref`), the cited facts by their paper names (C2, C3, C4;
% the ledger's A3, A5, A9; CHANGED Q-0281 wave 2: blind review W2F-3), and the blueprint-only split nodes (lem:gw-scheme-claims,
% lem:gw-kernels-self-decomposable, lem:stable-step-order, lem:cone-orders-psd, lem:stable-order-psd)
% replaced by the paper's clauses they copy. Nothing is shortened.
%   - prop:invariant-orders, lem:spectral-cones (statement moved here from Appendix B, where it served
%     only this proof), thm:cone-orders, prop:stable-order-generated: whole; their sketches in
%     Appendix B are deleted.
%   - prop:gw-ray-families: whole; its sketch in Appendix A is deleted.
%   - lem:stable-tail, prop:orbit-step-cone, thm:wedge-axis: whole; the main text kept its sketch,
%     with a pointer here, until the round-1 review (Gemini; the author's decision D3, 2026-11-03),
%     since when it states the idea of each proof in two or three sentences and points here.
%   - prop:similarity-stable-order-families: clauses (1) and (2), which thm:wedge-axis spends;
%     Appendix B keeps the sketch of (3)-(5). Since the round-2 review (D8, 2026-11-03) those are the
%     node prop:similarity-stable-order-kernels, and the proof here is of the whole proposition.
%   - prop:anisotropy-budget: clause (1), which thm:wedge-axis spends; Appendix B keeps clause (2).
% The other statements the proof spends were already proved completely where they stand
% (the list in the opening paragraph).

\section{The complete proof of the general linear row}\label{app:wedge-axis}

Section~\ref{sec:wedge-axis-proof} states Theorem~\ref{thm:wedge-axis} and the idea of its proof. This
appendix proves it completely, together with every statement the proof uses whose proof the main text
does not give in full, so that every step can be checked here. The proofs are in the order of their dependence.
\begin{itemize}
\item \S\ref{app:wedge-axis-orders}, invariant orders and their cones:
Proposition~\ref{prop:invariant-orders}, Lemma~\ref{lem:spectral-cones}, Theorem~\ref{thm:cone-orders}
and Proposition~\ref{prop:stable-order-generated}, which give the chain of cones and its orders.
\item \S\ref{app:wedge-axis-rays}, the ray column: Proposition~\ref{prop:gw-ray-families}, the base of
the row.
\item \S\ref{app:wedge-axis-steps}, the step set of an orbit: Lemma~\ref{lem:stable-tail},
Proposition~\ref{prop:orbit-step-cone},
Proposition~\ref{prop:similarity-stable-order-families} and clause (1) of
Proposition~\ref{prop:anisotropy-budget}, which decide which laws give families along each cone.
\item \S\ref{app:wedge-axis-theorem}, the theorem.
\end{itemize}
The other statements it uses are proved completely where they stand:
Theorems~\ref{thm:rank-deficient-step} and~\ref{thm:affine-rigidity} (\S\ref{sec:affine-single-step}),
Proposition~\ref{prop:stable-step-cone} (\S\ref{sec:affine-stable-cone}),
Lemma~\ref{lem:gw-kernels-from-origin}, Theorem~\ref{thm:cascade-infinitely-divisible} and the three
lemmas of the single step, Lemmas~\ref{lem:charfn-one-on-subspace}, \ref{lem:affine-reduction}
and~\ref{lem:single-ratio-cauchy} (Appendix~\ref{app:classical}), Lemma~\ref{lem:elliptical-stable} and
Proposition~\ref{prop:stable-exponent-cone} (Appendix~\ref{app:proofs}), and clause (1) of
Proposition~\ref{prop:tempered-step-cone}, whose short proof is in Appendix~\ref{app:orbits}.
% CHANGED Q-0281 wave 2: blind review W2F-3. Writing-rules pass (rule 4): the sentences naming the
% cited facts the proofs rest on and the standard tools they use are removed; a mid-line comment had
% broken them in print. Section 1.1 and section 2.6 hold that accounting.

\subsection{Invariant orders and their cones}\label{app:wedge-axis-orders}

This subsection proves that the invariant orders on $\PD(d)$ are read off their step sets, that the
orders generated by cone fields are those of the closed invariant cones, and that the chain $C_\alpha$
generates the stable orders and the Loewner order.

\begin{proof}[Proof of Proposition~\ref{prop:invariant-orders}]
\emph{(1).} Let $AA^{\tp} = \Sigma$ and $Q := \Sigma^{-1/2}A$; then $QQ^{\tp} =
\Sigma^{-1/2}AA^{\tp}\Sigma^{-1/2} = I$, so $Q \in \Orth(d)$ and $A^{-1}\Sigma'A^{-\tp} =
Q^{\tp}\Sigma^{-1/2}\Sigma'\Sigma^{-1/2}Q$, which lies in $\mathcal K$ exactly when the relative step does,
$\mathcal K$ being a step set. For $B \in \GL(d)$, $BA$ satisfies $(BA)(BA)^{\tp} = B\Sigma B^{\tp}$ and
$(BA)^{-1}B\Sigma'B^{\tp}(BA)^{-\tp} = A^{-1}\Sigma'A^{-\tp}$; with what was just shown, this is the
invariance. Let $R$ be a $\GL(d)$-invariant relation and $\mathcal K := \{M : I \,R\,M\}$. For
$Q \in \Orth(d)$, $QIQ^{\tp} = I$, so $I \,R\,M$ gives $I \,R\,QMQ^{\tp}$: $\mathcal K$ is a
step set. By invariance under $B = \Sigma^{-1/2}$, $\Sigma \,R\,\Sigma'$ if and only if
$I \,R\,\Sigma^{-1/2}\Sigma'\Sigma^{-1/2}$, that is, $\Sigma \le_{\mathcal K} \Sigma'$. A step set
is recovered from its relation as $\{M : I \le_{\mathcal K} M\}$, which gives uniqueness.

\emph{(2).} The relative step of $(\Sigma, \Sigma)$ is $I$. For antisymmetry let $M$ be the relative step
of $(\Sigma, \Sigma')$. The matrix $A := \Sigma^{1/2}M^{1/2}$ has $AA^{\tp} = \Sigma^{1/2}M\Sigma^{1/2} =
\Sigma'$ and $A^{-1}\Sigma A^{-\tp} = M^{-1/2}(\Sigma^{-1/2}\Sigma\Sigma^{-1/2})M^{-1/2} = M^{-1}$, so by
(1) $\Sigma' \le_{\mathcal K} \Sigma$ if and only if $M^{-1} \in \mathcal K$. The
relative step $M$ is $I$ exactly when $\Sigma = \Sigma'$, and every matrix of $\PD(d)$ is the relative step
of a pair, $N$ that of $(I, N)$. So $\le_{\mathcal K}$ is antisymmetric exactly when no $M \ne I$ has $M$
and $M^{-1}$ in $\mathcal K$. The map $\Sigma \mapsto \Sigma^{1/2}$ is continuous on $\PSD(d)$: if
$\Sigma_n \to \Sigma$, the entries of $\Sigma_n^{1/2}$ are bounded, $(\Sigma_n^{1/2})_{ij}^2 \le
(\Sigma_n)_{ii}$, and every limit point of $(\Sigma_n^{1/2})$ is positive semidefinite with square
$\Sigma$, hence is $\Sigma^{1/2}$ by the uniqueness of the positive square root. So
$(\Sigma, \Sigma') \mapsto \Sigma^{-1/2}\Sigma'\Sigma^{-1/2}$ is continuous on $\PD(d) \times \PD(d)$, and a closed
$\mathcal K$ gives a closed relation; conversely $\mathcal K$ is the slice $\{M : (I, M) \in {\le_{\mathcal K}}\}$
of the relation.

\emph{(3).} Put $A := \Sigma^{1/2}M^{1/2}$, so that $AA^{\tp} = \Sigma^{1/2}M\Sigma^{1/2} = \Sigma'$,
and, as in (1), $A = \Sigma'^{1/2}Q$ with $Q \in \Orth(d)$. Then $\Sigma'' = \Sigma'^{1/2}N\Sigma'^{1/2}
= AQ^{\tp}NQA^{\tp} = \Sigma^{1/2}M^{1/2}N'M^{1/2}\Sigma^{1/2}$, so the relative step of
$(\Sigma, \Sigma'')$ is $M^{1/2}N'M^{1/2} = M^{-1/2}(MN')M^{1/2}$, which has the eigenvalues of $MN'$.
The converse chain has relative steps $M$ and $M^{-1/2}(M^{1/2}NM^{1/2})M^{-1/2} = N$. So transitivity
is the stated closure property, $N' \in \mathcal K$ ranging over $\mathcal K$ as $N$ does. If it holds
and $A_1, A_2 \in \mathcal S_{\mathcal K}$, put $M_1 := A_1A_1^{\tp}$ and write $A_1 = M_1^{1/2}Q_1$ with
$Q_1 \in \Orth(d)$. Then $(A_1A_2)(A_1A_2)^{\tp} = M_1^{1/2}(Q_1A_2A_2^{\tp}Q_1^{\tp})M_1^{1/2}$ with
$Q_1A_2A_2^{\tp}Q_1^{\tp}$ in $\mathcal K$, so $A_1A_2 \in \mathcal S_{\mathcal K}$. Conversely, if $\mathcal S_{\mathcal K}$ is closed under
products and $M, N \in \mathcal K$, then $M^{1/2}, N^{1/2} \in \mathcal S_{\mathcal K}$, and so is
$M^{1/2}N^{1/2}$, that is, $M^{1/2}NM^{1/2} \in \mathcal K$. For the example,
$\diag(2,1)\diag(4,1)\diag(2,1) = \diag(16,1)$ and $\diag(2,1)\diag(1,4)\diag(2,1) = \diag(4,4)$.

\emph{(4).} $e^{tQXQ^{\tp}} = Qe^{tX}Q^{\tp}$ and $e^{t(cX)} = e^{(tc)X}$ give the two invariances. If
$X_n \to X$ with $X_n \in C(\mathcal K)$, then $e^{tX_n} \to e^{tX}$ in $\PD(d)$ for each $t$, so a closed
$\mathcal K$ contains $e^{tX}$. If $X$ and $-X$ lie in $C(\mathcal K)$, then $e^X$ and its inverse
$e^{-X}$ lie in $\mathcal K$, so antisymmetry gives $e^X = I$, hence $X = 0$, the exponential being
injective on $\Sym(d)$.
\end{proof}

The invariant cones are read through their eigenvalues. The lemma below gives the two matrix
inequalities the correspondence needs: a minimum principle for the partial sums of the eigenvalues of a
sum, and the majorization of the logarithmic spectrum of a product, which rests on the cited
fact C4.

% shared with blueprint lem:spectral-cones
\begin{lemma}[invariant cones, and the spectra of sums and products]\label{lem:spectral-cones}
Let $d \ge 1$. For $X \in \Sym(d)$ write $\lambda^\uparrow(X) = (\lambda^\uparrow_1(X), \dots,
\lambda^\uparrow_d(X))$ for its eigenvalues in increasing order, and for $M \in \PD(d)$ put
$\ell(M) := \lambda^\uparrow(\log M) = (\log\lambda^\uparrow_1(M), \dots, \log\lambda^\uparrow_d(M))$. For
$x \in \RR^d$ let $x^\uparrow$ and $x^\downarrow$ be its increasing and decreasing rearrangements,
$s_k(x) := x^\uparrow_1 + \dots + x^\uparrow_k$, and write $x \prec y$ when $s_k(x) \ge s_k(y)$ for
$1 \le k < d$ and $s_d(x) = s_d(y)$. A set $D \subseteq \RR^d$ is \emph{symmetric} if it is invariant
under permutations of the coordinates, and $D^\# := \{w \in \RR^d : \langle w, x\rangle \ge 0 \text{ for
every } x \in D\}$.
\begin{enumerate}
\item \emph{(Sums.)} $s_k(\lambda^\uparrow(X)) = \min\{\tr(U^{\tp}XU) : U \in \RR^{d\times k},\
U^{\tp}U = I\}$. Consequently $|s_k(\lambda^\uparrow(X)) - s_k(\lambda^\uparrow(Y))| \le k\|X - Y\|$, with
$\|\cdot\|$ the operator norm, and $\lambda^\uparrow(X + Y) \prec \lambda^\uparrow(X) + \lambda^\uparrow(Y)$.
\item \emph{(Products.)} For $M, N \in \PD(d)$, $\ell(M^{1/2}NM^{1/2}) \prec \ell(M) + \ell(N)$.
\item \emph{(Symmetric cones are closed under majorization.)} Let $D \subseteq \RR^d$ be a symmetric
closed convex cone. Then $x \in D$ if and only if $\sum_iw^\downarrow_ix^\uparrow_i \ge 0$ for every
$w \in D^\#$. Consequently $y \in D$ and $x \prec y$ imply $x \in D$.
\item \emph{(Cones in $\Sym(d)$ and in eigenvalue space.)} For a symmetric closed convex cone
$D \subseteq \RR^d$, $C_D := \{X \in \Sym(d) : \lambda^\uparrow(X) \in D\}$ is a closed convex cone
invariant under $X \mapsto QXQ^{\tp}$, $Q \in \Orth(d)$. Every closed convex cone $C \subseteq \Sym(d)$
with that invariance is $C_D$ for exactly one such $D$, namely $D = \{x : \diag(x) \in C\}$, and
$C \cap (-C) = \{0\}$ if and only if $D \cap (-D) = \{0\}$. If moreover $C \ne \{0\}$ and
$C \cap (-C) = \{0\}$, then exactly one of $I$ and $-I$ lies in $C$, and $\tr X$ has that sign for
every $X \in C \setminus \{0\}$.
\item \emph{(Exponentials.)} For $C = C_D$ as in (4), $\exp C := \{e^X : X \in C\}$ equals
$\{M \in \PD(d) : \ell(M) \in D\}$; it is a step set (Proposition~\ref{prop:invariant-orders}), closed in $\PD(d)$,
and $M^{1/2}NM^{1/2} \in \exp C$ whenever $M, N \in \exp C$.
\end{enumerate}
\end{lemma}

Clause (5) is the transitivity of every order whose steps are the exponentials of an invariant closed
convex cone (Proposition~\ref{prop:invariant-orders}(3)). Clause (4) reads such a cone through the
permutation-symmetric cone of its eigenvalues. Its last sentence says that every invariant closed
convex cone $C \ne \{0\}$ with $C \cap (-C) = \{0\}$ contains $I$ or $-I$, the directions of
$\Sym(d)$ fixed by $\Orth(d)$. Clause (1) is Ky Fan's extremal principle with his majorization of
the eigenvalues of a sum (\sscite{marshall2011inequalities}, 20.A.2, p.~785, and 9.G.1, p.~329).
Clause (2) is the log-majorization of the eigenvalues of a product (9.H.1.a, pp.~338--339). The
correspondence in (4) follows from Davis's theorem on orthogonally invariant convex functions
(\sscite{davis1957all}, Cors.~1--2, p.~277).

\begin{proof}[Proof of Lemma~\ref{lem:spectral-cones}]
\emph{(1).} Write $X = \sum_j\lambda_jv_jv_j^{\tp}$ with $\lambda_j = \lambda^\uparrow_j(X)$ and $(v_j)$ an
orthonormal basis. For $U$ with $U^{\tp}U = I$, $\tr(U^{\tp}XU) = \sum_j\lambda_jp_j$ with
$p_j := |U^{\tp}v_j|^2$; since $UU^{\tp}$ is an orthogonal projection, $0 \le p_j \le 1$, and
$\sum_jp_j = \tr(U^{\tp}U) = k$. Then
\[
  \sum_j\lambda_jp_j - \sum_{j\le k}\lambda_j = \sum_{j\le k}\lambda_j(p_j - 1) + \sum_{j>k}\lambda_jp_j
  \ge \lambda_k\sum_{j\le k}(p_j - 1) + \lambda_k\sum_{j>k}p_j = 0 ,
\]
using $p_j - 1 \le 0$ and $\lambda_j \le \lambda_k$ for $j \le k$, and $\lambda_j \ge \lambda_k$ for
$j > k$; equality holds for $U = (v_1, \dots, v_k)$. Each $X \mapsto \tr(U^{\tp}XU)$ is linear and
$k\|\cdot\|$-Lipschitz, so their minimum is $k\|\cdot\|$-Lipschitz, and the minimum of a sum is at least
the sum of the minima: $s_k(\lambda^\uparrow(X+Y)) \ge s_k(\lambda^\uparrow(X)) +
s_k(\lambda^\uparrow(Y))$, with equality at $k = d$, where both sides are traces. The vector
$\lambda^\uparrow(X) + \lambda^\uparrow(Y)$ is increasing, so its $s_k$ is the right-hand side.

\emph{(2).} Let $P := M^{1/2}NM^{1/2}$. The products $MN = M^{1/2}(M^{1/2}N)$ and $(M^{1/2}N)M^{1/2} = P$
have the same characteristic polynomial, so the eigenvalues of $MN$, with multiplicity, are those of $P$.
By the inequality for the eigenvalues of a product in its weak form (\ledger{A9}), for $1 \le k \le d$
\[
  \sum_{i\le k}\log\lambda^\downarrow_i(P) \le \sum_{i\le k}\bigl(\log\lambda^\downarrow_i(M) +
  \log\lambda^\downarrow_i(N)\bigr) ,
\]
with $\lambda^\downarrow$ the eigenvalues in decreasing order: the sum of the $k$ largest entries of
$\ell(P)$ is at most that of $\ell(M) + \ell(N)$, both vectors being increasing. The totals are equal,
$\sum_i\ell_i(P) = \log\det P = \log(\det M\det N) = \sum_i(\ell_i(M) + \ell_i(N))$. Subtracting from the
totals, $s_j(\ell(P)) \ge s_j(\ell(M) + \ell(N))$ for $j = d - k$, $1 \le j < d$, with equality at $j = d$,
which is the claim.

\emph{(3).} \emph{The dual description.} If $x \in D$ and $w \in D^\#$, every rearrangement of $x$ lies in
$D$, so $\langle w, \sigma x\rangle \ge 0$ for every permutation $\sigma$; the number
$\sum_iw^\downarrow_ix^\uparrow_i$ is one of these, the largest entries of $w$ meeting the smallest of $x$,
and it is the least of them (the rearrangement inequality). Conversely let $x \notin D$. As $D$ is closed
and convex, a linear functional $f$ separates $x$ from $D$ strictly: $f(x) < u < f(z)$ for a number $u$ and
every $z \in D$ (Hahn--Banach). Write $f(z) = \langle w, z\rangle$. With $z = 0$, $u < 0$; and if
$\langle w, z\rangle < 0$ for some $z \in D$, then $cz \in D$ with $c := u/\langle w, z\rangle > 0$ gives
$u < \langle w, cz\rangle = u$, which is absurd. So $w \in D^\#$ and $\langle w, x\rangle < u < 0$, and by
the rearrangement inequality $\sum_iw^\downarrow_ix^\uparrow_i \le \langle w, x\rangle < 0$.
\emph{Majorization.} By Abel summation, for $w \in \RR^d$ and $x \in \RR^d$,
\[
  \sum_{i=1}^dw^\downarrow_ix^\uparrow_i = \sum_{k=1}^{d-1}\bigl(w^\downarrow_k - w^\downarrow_{k+1}\bigr)s_k(x)
  + w^\downarrow_d\,s_d(x) ,
\]
with every coefficient $w^\downarrow_k - w^\downarrow_{k+1} \ge 0$. If $x \prec y$, each term for $x$ is at
least the corresponding term for $y$ (the last ones are equal), so
$\sum_iw^\downarrow_ix^\uparrow_i \ge \sum_iw^\downarrow_iy^\uparrow_i \ge 0$ for $w \in D^\#$ when
$y \in D$, and $x \in D$ by the dual description.

\emph{(4).} $\lambda^\uparrow(QXQ^{\tp}) = \lambda^\uparrow(X)$ and $\lambda^\uparrow(cX) = c\lambda^\uparrow(X)$
for $c \ge 0$. Each $\lambda^\uparrow_k = s_k - s_{k-1}$ is continuous by (1), and $D$ is closed, so $C_D$
is closed. If $X, Y \in C_D$, then $\lambda^\uparrow(X) + \lambda^\uparrow(Y) \in D$, $D$ being a convex
cone, and (1) with (3) gives $\lambda^\uparrow(X+Y) \in D$. So $C_D$ is a convex cone. Given $C$, the set
$D := \{x : \diag(x) \in C\}$ is a closed convex cone, symmetric because permutation matrices are
orthogonal. Writing $X = Q\diag(\lambda^\uparrow(X))Q^{\tp}$, $X \in C$ if and only if $\lambda^\uparrow(X)
\in D$, so $C = C_D$. And $D$ is determined by $C_D$ as $\{x : \diag(x) \in C_D\}$. If $0 \ne x \in D \cap
(-D)$ then $\diag(x) \in C \cap (-C)$; if $0 \ne X \in C \cap (-C)$ then $\lambda^\uparrow(X)$ and
$\lambda^\uparrow(-X)$, a rearrangement of $-\lambda^\uparrow(X)$, lie in $D$, so $0 \ne \lambda^\uparrow(X)
\in D \cap (-D)$. For the last sentence let $0 \ne X \in C$ and $x := \lambda^\uparrow(X) \ne 0$. For
$m = 0, \dots, d-1$ let $\sigma_mx$ be the cyclic shift $(\sigma_mx)_i := x_{i+m}$, indices modulo $d$.
Each $\sigma_mx$ lies in $D$, and their sum is $(\tr X)\mathbf 1$, with $\mathbf 1 = (1, \dots, 1)$, since
every coordinate of it is $\sum_jx_j$; so $(\tr X)\mathbf 1 \in D$. If $\tr X = 0$ this sum is $0$, so
$x = \sigma_0x = -\sum_{m\ne0}\sigma_mx \in -D$ and $x \in D \cap (-D)$, which is excluded. So $\tr X \ne 0$,
and $(\operatorname{sign}\tr X)\mathbf 1 \in D$, that is, $(\operatorname{sign}\tr X)\,I \in C$. As $I$ and
$-I$ cannot both lie in $C$, the sign is the same for every nonzero $X \in C$, which is the claim.

\emph{(5).} For $X \in \Sym(d)$, $e^X$ has the eigenvalues $e^{\lambda_i(X)}$, and $\log$ is the inverse of
$\exp$ from $\Sym(d)$ onto $\PD(d)$; so $M \in \exp C$ if and only if $\log M \in C$, that is,
$\ell(M) \in D$. Orthogonal conjugation commutes with $\exp$, so $\exp C$ is a step set. On $\PD(d)$,
$\ell(M) = (\log\lambda^\uparrow_i(M))_i$ is continuous, each $\lambda^\uparrow_i$ being continuous by (1)
and positive there, and $D$ is closed, so $\exp C = \{M : \ell(M) \in D\}$ is closed in $\PD(d)$. For
$M, N \in \exp C$, $\ell(M) + \ell(N) \in D$, and (2) with (3) gives $\ell(M^{1/2}NM^{1/2}) \in D$.
\end{proof}

\begin{proof}[Proof of Theorem~\ref{thm:cone-orders}]
Write $C = C_D$ (Lemma~\ref{lem:spectral-cones}(4)) and, for $w \in D^\#$ and $\Sigma, \Sigma' \in \PD(d)$,
\[
  F_w(\Sigma, \Sigma') := \sum_{i=1}^dw^\downarrow_i\,\ell_i\bigl(\Sigma^{-1/2}\Sigma'\Sigma^{-1/2}\bigr)
  = \sum_{k=1}^{d-1}\bigl(w^\downarrow_k - w^\downarrow_{k+1}\bigr)s_k(\ell) + w^\downarrow_d\,s_d(\ell) ,
\]
$\ell$ the logarithmic spectrum of the relative step, increasing. By Lemma~\ref{lem:spectral-cones}(5) and (3),
$\Sigma \le_{\exp C} \Sigma'$ if and only if $F_w(\Sigma, \Sigma') \ge 0$ for every $w \in D^\#$.
\emph{Superadditivity.} For any $\Sigma, \Sigma', \Sigma'' \in \PD(d)$, with $M$, $N$ the relative steps of
$(\Sigma, \Sigma')$, $(\Sigma', \Sigma'')$, the relative step of $(\Sigma, \Sigma'')$ is
$M^{1/2}N'M^{1/2}$ with $N'$ orthogonally conjugate to $N$ (the computation of
Proposition~\ref{prop:invariant-orders}(3), which does not use the order); by Lemma~\ref{lem:spectral-cones}(2) its $s_k$ is
at least $s_k(\ell(M)) + s_k(\ell(N))$ for $k < d$, with equality at $k = d$. The coefficients
$w^\downarrow_k - w^\downarrow_{k+1}$ are nonnegative, so
$F_w(\Sigma, \Sigma'') \ge F_w(\Sigma, \Sigma') + F_w(\Sigma', \Sigma'')$.

\emph{(1), the geodesic.} Let $M \in \exp C$, $X := \log M \in C$ and $\gamma(t) = \Sigma^{1/2}e^{tX}\Sigma^{1/2}$.
With $A := \Sigma^{1/2}e^{tX/2}$, $\gamma(t) = AA^{\tp}$ and $\gamma'(t) = \Sigma^{1/2}Xe^{tX}\Sigma^{1/2} =
AXA^{\tp}$, so $\gamma(t)^{-1/2}\gamma'(t)\gamma(t)^{-1/2} = QXQ^{\tp}$ with $Q := \gamma(t)^{-1/2}A$, and
$QQ^{\tp} = \gamma(t)^{-1/2}AA^{\tp}\gamma(t)^{-1/2} = I$; so it lies in $C$. The entries of $\gamma$ are, in an
eigenbasis of $X$, finite sums of exponentials $e^{tx_k}$, so $\gamma$ is smooth.
\emph{(1), causal curves end in the order.} Let $\gamma$ be $C$-causal and $w \in D^\#$, and put
$g(t) := F_w(\gamma(0), \gamma(t))$, so $g(0) = 0$. On $[0,1]$, $g(t) = \sum_iw^\downarrow_i\log
\lambda^\uparrow_i(N(t))$ with $N(t) := \gamma(0)^{-1/2}\gamma(t)\gamma(0)^{-1/2}$. The least eigenvalue of
$N(t)$ is continuous in $t$ (Lemma~\ref{lem:spectral-cones}(1)) and positive, so it is at least some $m > 0$ on
$[0,1]$, and $\log$ is $(1/m)$-Lipschitz on $[m, \infty)$. Each $\lambda^\uparrow_i$ is Lipschitz in the
operator norm by Lemma~\ref{lem:spectral-cones}(1), and the operator norm of $N(s) - N(t)$ is at most
$\|\gamma(0)^{-1/2}\|^2\sum_{a,b}|\gamma_{ab}(s) - \gamma_{ab}(t)|$. So $|g(s) - g(t)|$ is bounded by a
constant times $\sum_{a,b}|\gamma_{ab}(s) - \gamma_{ab}(t)|$, and $g$ is absolutely continuous because the
entries of $\gamma$ are. Let $t \in (0,1)$ be a point at which the entries of $\gamma$ are differentiable and
$Y := \gamma(t)^{-1/2}\gamma'(t)\gamma(t)^{-1/2} \in C$; almost every $t$ is such. For $h > 0$,
$N_h := \gamma(t)^{-1/2}\gamma(t+h)\gamma(t)^{-1/2} = I + hY + o(h)$, and as $\lambda^\uparrow(I + hY) =
\mathbf 1 + h\lambda^\uparrow(Y)$, the Lipschitz bound of Lemma~\ref{lem:spectral-cones}(1) gives
$x_i(h) := \lambda^\uparrow_i(N_h) = 1 + h\lambda^\uparrow_i(Y) + o(h)$. From $(x - 1)/x \le \log x \le x - 1$
for $x > 0$, $\log x_i(h) = h\lambda^\uparrow_i(Y) + o(h)$. Superadditivity gives
$g(t+h) - g(t) \ge F_w(\gamma(t), \gamma(t+h))$, hence, where $g$ is differentiable at $t$,
\[
  g'(t) \ge \lim_{h\downarrow0}\frac1hF_w(\gamma(t), \gamma(t+h)) = \sum_iw^\downarrow_i\lambda^\uparrow_i(Y) \ge 0 ,
\]
by Lemma~\ref{lem:spectral-cones}(3), as $\lambda^\uparrow(Y) \in D$. An absolutely continuous function is the
integral of its derivative, which exists almost everywhere; so $g(1) = \int_0^1g' \ge 0$, and as $w$ was
arbitrary, $\gamma(0) \le_{\exp C} \gamma(1)$. \emph{(1), the rest.} $\exp C$ is a closed step set
containing $I$ (Lemma~\ref{lem:spectral-cones}(5)), so $\le_{\exp C}$ is closed, reflexive and, by
Lemma~\ref{lem:spectral-cones}(5) and Proposition~\ref{prop:invariant-orders}(3), transitive; $M$ and $M^{-1}$ lie in
$\exp C$ exactly when $\log M \in C \cap (-C)$, which with Proposition~\ref{prop:invariant-orders}(2) is the criterion
for antisymmetry; and $e^{tX} \in \exp C$ if and only if $tX \in C$, $\log$ being injective, so
$C(\exp C) = C$.

\emph{(2).} Let $\mathcal L$ be the set of $B$ named. $\mathcal S_{\mathcal K}$ is closed in $\GL(d)$, as
the preimage of the closed set $\mathcal K$ under $A \mapsto AA^{\tp}$; it contains $\Orth(d)$, as
$I \in \mathcal K$; and it is closed under products by Proposition~\ref{prop:invariant-orders}(3). $\mathcal L$ is closed
under positive multiples and, $\mathcal S_{\mathcal K}$ being closed, under limits. It is closed under
sums by the Lie product formula $(e^{B/n}e^{B'/n})^n \to e^{B+B'}$. Put $P_n := e^{B/n}e^{B'/n}$,
$Q_n := e^{(B+B')/n}$ and $a := \|B\| + \|B'\|$. The curve $\varphi(s) := e^{sB}e^{sB'} - e^{s(B+B')}$ has
$\varphi(0) = 0$ and $\varphi'(0) = B + B' - (B + B') = 0$, so $n\|P_n - Q_n\| = n\|\varphi(1/n)\| \to 0$; and
$\|P_n\|, \|Q_n\| \le e^{a/n}$, as $\|e^Z\| \le e^{\|Z\|}$. From $P^n - Q^n = P(P^{n-1} - Q^{n-1}) + (P - Q)Q^{n-1}$,
by induction,
\[
  \|P_n^n - Q_n^n\| \le n\,(e^{a/n})^n\,\|P_n - Q_n\| = e^a\,n\|P_n - Q_n\| \to 0 ,
\]
and $Q_n^n = e^{B+B'}$. Applied to $tB$ and $tB'$, the products $(e^{tB/n}e^{tB'/n})^n$ lie in $\mathcal S_{\mathcal K}$ and tend to $e^{t(B+B')}$. Three cases then describe $\mathcal L$.
\begin{itemize}
\item For symmetric $X$, $e^{tX}(e^{tX})^{\tp} = e^{2tX}$, so $X \in \mathcal L$ exactly when
$X \in C(\mathcal K)$.
\item For skew $K$, $e^{tK} \in \Orth(d) \subseteq \mathcal S_{\mathcal K}$, so $\pm K \in \mathcal L$.
\item Adding the skew part of $B$ or its negative, $B \in \mathcal L$ if and only if
$\frac12(B + B^{\tp}) \in \mathcal L$.
\end{itemize} So $\mathcal L$ is as stated, and $C(\mathcal K) = \mathcal L \cap
\Sym(d)$ is a convex cone; it is closed, invariant and pointed by Proposition~\ref{prop:invariant-orders}(4). By
(1), $\preceq_{C(\mathcal K)}$ is $\le_{\exp C(\mathcal K)}$, and $\exp C(\mathcal K) \subseteq \mathcal K$ by
definition. (c)$\Rightarrow$(b): let $M \in \mathcal K$ and $X := \log M$. Then $e^{tX} = M^t$ lies in
$\mathcal K$ for $t \in [0,1]$ by (c). For $t = j + \tau$ with $j \ge 1$ an integer and $0 \le \tau < 1$,
induction on $j$ gives it: $M^{j+\tau} = M^{1/2}M^{j-1+\tau}M^{1/2}$ is the composite of the steps $M$ and
$M^{j-1+\tau}$ of $\mathcal K$, which lies in $\mathcal K$ by transitivity
(Proposition~\ref{prop:invariant-orders}(3)). So $X \in C(\mathcal K)$ and $\mathcal K \subseteq \exp C(\mathcal K)$. (b)$\Rightarrow$(a) is (1) with $C' = C(\mathcal K)$.
(a)$\Rightarrow$(c): by (1), $\mathcal K = \exp C'$, closed under $M \mapsto M^\tau$ as $C'$ is a cone, and
$C' = C(\exp C') = C(\mathcal K)$.

\emph{(3).} For $M, N$ with least eigenvalues at least $2$ and a unit $x$,
$\langle x, M^{1/2}NM^{1/2}x\rangle \ge 2|M^{1/2}x|^2 \ge 4$, so composites stay in $\mathcal K$, and
composites with $I$ trivially do: $\le_{\mathcal K}$ is transitive (Proposition~\ref{prop:invariant-orders}(3)), and
reflexive. $\mathcal K$ is a step set, and closed since $\lambda^\uparrow_1$ is continuous. If $M \ne I$ and
$M \in \mathcal K$, then $\lambda^\uparrow_1(M^{-1}) = 1/\lambda^\uparrow_d(M) \le 1/2$, so $M^{-1} \notin
\mathcal K$: antisymmetry. For $X \ne 0$ and small $t > 0$, $e^{tX} \ne I$ and its least eigenvalue
$e^{t\lambda^\uparrow_1(X)}$ is below $2$, so $X \notin C(\mathcal K)$. A $\{0\}$-causal curve has
derivative $0$ almost everywhere, and an absolutely continuous curve with derivative $0$ almost everywhere is
constant, so $\preceq_{\{0\}}$ is equality. If $\le_{\mathcal K}$ were $\preceq_{C'}$, then for $X \in C'$ and
$t \ge 0$ the geodesic of (1) from $I$ to $e^{tX}$ would be $C'$-causal, so $I \le_{\mathcal K} e^{tX}$, that is
$e^{tX} \in \mathcal K$: $C' \subseteq C(\mathcal K) = \{0\}$. Then $\le_{\mathcal K}$ would be equality,
while $I \le_{\mathcal K} 2I$.

\emph{(4).} By Lemma~\ref{lem:spectral-cones}(4) every nonzero $X \in C$ has $\tr X > 0$.
On $\PD(d)$ the second case is contained in the third, since $\Sigma' = c\Sigma$ has relative step
$cI = e^{(\log c)I}$ with $(\log c)I \in C$. If $\Sigma$ is definite, so is every $\Sigma'$ above it. If
$\Sigma \ne 0$ is singular, the $\Sigma'$ above it are its multiples $c\Sigma$, $c \ge 1$. The claims follow.
\begin{itemize}
\item Reflexivity is $c = 1$.
\item Transitivity: from $0$ trivially, from a nonzero singular matrix along its multiples, and on
$\PD(d)$ by that of $\preceq_C$.
\item Antisymmetry: on $\PD(d)$ it is (1). $\Sigma' \preceq^+_C 0$ forces $\Sigma' = 0$. And $c\Sigma = \Sigma$
with $\Sigma = c'(c\Sigma)$, $c, c' \ge 1$, forces $c = 1$.
\item Invariance: $A \cdot 0 = 0$, $A \cdot (c\Sigma) = c\,A\cdot\Sigma$, and $\preceq_C$ is invariant by (1).
$G$-invariance follows from $\GL(d)$-invariance.
\item The rays: $c\Sigma \preceq^+_C c'\Sigma$ for $0 \le c \le c'$, by the first case if $c = 0$ and the
second otherwise.
\item Chains: $s\Sigma'$ from $0$ and $(1 + s(c-1))\Sigma$ along a ray are chains by
Theorem~\ref{thm:cascade-infinitely-divisible}(3)(b). On $\PD(d)$ the geodesic $\gamma$ of (1) is one, by
Theorem~\ref{thm:cascade-infinitely-divisible}(3)(c) applied to $\mathcal K = \exp C$, a step set whose order is a partial
order by (1) and which is closed under $M \mapsto M^\tau$, $0 \le \tau \le 1$, as $C$ is a cone; and
$\le_{\exp C}$ is $\preceq_C$ by (1). So (P) holds.
\item The Loewner order: if $C \subseteq \PSD(d)$, then $e^X \ge I$ for $X \in C$, so every pair of
$\preceq^+_C$ is a Loewner pair. If $X \in C$ is not positive semidefinite, $I \preceq^+_C e^X$ while
$e^X - I$ is not positive semidefinite.
\item The ray order: for $C = \{cI : c \ge 0\}$, $\exp C = \{cI : c \ge 1\}$, so on $\PD(d)$ the third case
is the second, and $\preceq^+_C$ is the ray order.
\end{itemize}
\end{proof}

\begin{proof}[Proof of Proposition~\ref{prop:stable-order-generated}]
\emph{The cones.} The conditions defining $C_\alpha$ are invariant under orthogonal conjugation and positive
scaling; $\lambda_{\min}$ is concave and continuous and $\tr$ is linear, so $C_\alpha$ is a closed convex
cone; $I$ satisfies $d + \alpha \ge d$; and $C_\alpha \subseteq \PSD(d)$ meets $-\PSD(d)$ only in $0$.

\emph{Three facts about curves.} (a) Let $\eta : [0,1] \to \Sym(d)$ have absolutely continuous entries. If
$\eta'(t)$ is positive semidefinite for almost every $t$, then for each $x$ the function $\langle x,
\eta(t)x\rangle$ is absolutely continuous with derivative $\langle x, \eta'(t)x\rangle \ge 0$ almost
everywhere, so it is the integral of its derivative and $\eta(0) \le \eta(1)$. If $\eta' = 0$ almost
everywhere, each entry is constant, and $\eta(0) = \eta(1)$. (b) For $s \in \RR$ and a $C$-causal $\gamma$
($C$ any cone), the entries of $(\det\gamma)^s\gamma$ are absolutely continuous on $[0,1]$: $\det\gamma$ is a
polynomial in the entries of $\gamma$, positive and continuous, hence between two positive numbers on
$[0,1]$, where $x \mapsto x^s$ is Lipschitz. At almost every $t$ the entries of $\gamma$ are differentiable;
there, writing $\gamma(u) = \gamma(t)\cdot\gamma(t)^{-1}\gamma(u)$ and expanding the determinant of
$\gamma(t)^{-1}\gamma(u)$, which is $I$ at $u = t$, gives Jacobi's formula $(\det\gamma)'(t) =
\det\gamma(t)\,\tr(\gamma(t)^{-1}\gamma'(t))$. With $Y := \gamma^{-1/2}\gamma'\gamma^{-1/2}$, so that
$\gamma' = \gamma^{1/2}Y\gamma^{1/2}$ and $\tr(\gamma^{-1}\gamma') = \tr Y$,
\[
  \bigl((\det\gamma)^s\gamma\bigr)' = (\det\gamma)^s\bigl(\gamma' + s(\tr Y)\gamma\bigr)
  = (\det\gamma)^s\,\gamma^{1/2}\bigl(Y + s(\tr Y)I\bigr)\gamma^{1/2} .
\]
(c) If $C$ is invariant under orthogonal conjugation, $X \in C$ and $M = e^X$ is the relative step of
$(\Sigma, \Sigma')$, then $\gamma(t) := \Sigma^{1/2}M^t\Sigma^{1/2} = \Sigma^{1/2}e^{tX}\Sigma^{1/2}$ is a
$C$-causal curve from $\Sigma$ to $\Sigma'$: its entries are smooth, and with $A := \Sigma^{1/2}e^{tX/2}$,
$\gamma(t) = AA^{\tp}$ and $\gamma'(t) = AXA^{\tp}$, so $\gamma(t)^{-1/2}\gamma'(t)\gamma(t)^{-1/2} =
QXQ^{\tp}$ with $Q := \gamma(t)^{-1/2}A$ orthogonal.

\emph{(1).} $X \in C_0$ means $X \ge 0$ and $d\lambda_{\min}(X) \ge \tr X$, that is, every eigenvalue equals the
least one and is nonnegative: $X = cI$, $c \ge 0$. Let $\gamma$ be $C_0$-causal from $\Sigma$ to $\Sigma'$.
Almost everywhere $Y = cI$ for some $c \ge 0$, so $Y + s(\tr Y)I = 0$ for $s = -1/d$, and by (b) and (a)
$(\det\gamma)^{-1/d}\gamma$ is constant: $\Sigma' = c\,\Sigma$ with $c = (\det\Sigma'/\det\Sigma)^{1/d}$. As
$C_0 \subseteq \PSD(d)$, $\gamma' = \gamma^{1/2}Y\gamma^{1/2}$ is positive semidefinite almost everywhere, so
$\Sigma \le \Sigma'$ by (a), and $(c - 1)\Sigma \ge 0$ forces $c \ge 1$. Conversely the relative step of
$(\Sigma, c\Sigma)$ is $cI = e^{(\log c)I}$ with $(\log c)I \in C_0$ for $c \ge 1$, and (c) gives the curve.
For $0 < \alpha < \infty$, $\tr X \ge d\lambda_{\min}(X)$, so
$(d+\alpha)\lambda_{\min}(X) \ge \tr X$ gives $\alpha\lambda_{\min}(X) \ge 0$, hence $X \ge 0$. With
$M = e^X$, $\lambda_{\min}(M) = e^{\lambda_{\min}(X)}$ and $\det M = e^{\tr X}$, so $M \in \exp C_\alpha$ if and
only if $\lambda_{\min}(M)^{d+\alpha} \ge \det M$, which is condition (c) of Proposition~\ref{prop:stable-step-cone}(1)
when $0 < \alpha < 2$. Let $\gamma$ be $C_\alpha$-causal and $s := -1/(d+\alpha)$. Almost everywhere $Y \in
C_\alpha$, so $Y + s(\tr Y)I = Y - \frac{\tr Y}{d+\alpha}I \ge 0$, and by (b) and (a) $\Phi_\alpha(\Sigma) \le
\Phi_\alpha(\Sigma')$. Dividing by $(\det\Sigma)^s > 0$ and congruence by $\Sigma^{-1/2}$, $\Phi_\alpha(\Sigma)
\le \Phi_\alpha(\Sigma')$ is $I \le (\det M)^sM$, that is $\lambda_{\min}(M)(\det M)^s \ge 1$, which is
$\lambda_{\min}(M)^{d+\alpha} \ge \det M$. Conversely, if that holds, $M \in \exp C_\alpha$ and (c) gives the
curve. Finally $e^X \ge I$ if and only if $X \ge 0$, so $\exp C_\infty = \{M : M \ge I\}$, and
$M \ge I$ is $\Sigma \le \Sigma'$, congruence preserving positive semidefiniteness. A $\PSD(d)$-causal curve has
$\gamma' = \gamma^{1/2}Y\gamma^{1/2} \ge 0$ almost everywhere, so $\Sigma \le \Sigma'$ by (a); conversely
$\Sigma \le \Sigma'$ puts $M$ in $\exp C_\infty$, and (c) gives the curve.

\emph{(2).} Let $s := -1/(d+\alpha)$ and, for a matrix $B$, $B' := B + s(\tr B)I$. As
$|\det e^{tB}|^s = e^{st\tr B}$ and $I$ commutes with $B$, $\psi_\alpha(e^{tB}) = e^{tB'}$. By (1) and (3),
$e^{tB}e^{tB^{\tp}} \in \exp C_\alpha$ if and only if $I \preceq_{C_\alpha} e^{tB}(e^{tB})^{\tp}$, if and only if
$I \le \Phi_\alpha(e^{tB}(e^{tB})^{\tp}) = \psi_\alpha(e^{tB})\psi_\alpha(e^{tB})^{\tp} = e^{tB'}e^{tB'^{\tp}}$.
For $x \in \RR^d$ put $f(t) := |e^{tB'^{\tp}}x|^2 = \langle x, e^{tB'}e^{tB'^{\tp}}x\rangle$, whose derivative
is $\langle y, (B' + B'^{\tp})y\rangle$ with $y = e^{tB'^{\tp}}x$. If $\lambda_{\min}(\frac12(B' + B'^{\tp}))
\ge 0$ then $f(t) \ge f(0) = |x|^2$ for $t \ge 0$; conversely $f(t) \ge f(0)$ for $t \ge 0$ gives
$f'(0) = \langle x, (B' + B'^{\tp})x\rangle \ge 0$. So $e^{tB}e^{tB^{\tp}} \in \exp C_\alpha$ for every
$t \ge 0$ if and only if $\lambda_{\min}(\frac12(B + B^{\tp})) + s\tr B \ge 0$, that is
$(d+\alpha)\lambda_{\min}(\frac12(B + B^{\tp})) \ge \tr B$, which is \ref{eq:stable-exponent-cone}. A symmetric
$X$ has $\frac12(X + X^{\tp}) = X$, and the condition $X \ge 0$ is implied for $\alpha > 0$, so
$\mathfrak k_\alpha \cap \Sym(d) = C_\alpha$.

\emph{(3).} Put $s := -1/(d+\alpha)$, so $1 + ds = \alpha/(d+\alpha) > 0$. For $T := \Phi_\alpha(\Sigma)$,
$\det T = (\det\Sigma)^{1+ds}$, so $(\det T)^{1/\alpha} = (\det\Sigma)^{1/(d+\alpha)}$ and
$(\det T)^{1/\alpha}T = \Sigma$; in the same way $\Phi_\alpha((\det T)^{1/\alpha}T) = T$. Both maps are
smooth. By (1), $\Sigma \preceq_{C_\alpha} \Sigma'$ if and only if $\Phi_\alpha(\Sigma) \le \Phi_\alpha(\Sigma')$.
By Jacobi's formula of (b) applied to $\varepsilon \mapsto \Sigma + \varepsilon V$, the derivative of
$\Phi_\alpha$ at $\Sigma$ in the direction $V$ is
\[
  D\Phi_\alpha(\Sigma)V = (\det\Sigma)^s\bigl(V + s\tr(\Sigma^{-1}V)\,\Sigma\bigr) ,
  \qquad
  D\Phi_\alpha(\Sigma)\bigl(\Sigma^{1/2}Y\Sigma^{1/2}\bigr) = (\det\Sigma)^s\,\Sigma^{1/2}\bigl(Y + s(\tr Y)I\bigr)\Sigma^{1/2} .
\]
The second is positive semidefinite when $\lambda_{\min}(Y) \ge (\tr Y)/(d+\alpha)$, in particular for
$Y \in C_\alpha$. Conversely, for $P \ge 0$ put $Z := (\det\Sigma)^{-s}\Sigma^{-1/2}P\Sigma^{-1/2} \ge 0$ and
$Y := Z - \frac{s}{1+ds}(\tr Z)I$; then $\tr Y = \tr Z/(1+ds)$ and $Y + s(\tr Y)I = Z \ge 0$, so $Y \in
C_\alpha$, and $D\Phi_\alpha(\Sigma)(\Sigma^{1/2}Y\Sigma^{1/2}) = P$: every positive semidefinite direction is
attained.
For $A \in \GL(d)$, $\det(A\Sigma A^{\tp}) = (\det A)^2\det\Sigma$ gives the covariance. Finally
$\psi_\alpha(AB) = \psi_\alpha(A)\psi_\alpha(B)$, $|\det A|^s$ being a character and scalars central, and with
$s' := -s/(1+ds)$ the map $A \mapsto |\det A|^{s'}A$ is its inverse, since $|\det(|\det A|^{s'}A)|^s|\det A|^{s'} =
|\det A|^{s'(1+ds) + s} = 1$, and symmetrically for the other composition. As $\Phi_\alpha(I) = I$, $I \preceq_{C_\alpha} AA^{\tp}$ if and only if
$I \le \Phi_\alpha(AA^{\tp}) = \psi_\alpha(A)\psi_\alpha(A)^{\tp}$, which is the last sentence.
\end{proof}

\subsection{The ray column}\label{app:wedge-axis-rays}

The base of the general linear row is the ray order. Its families are the orbits of the
rotation-invariant self-decomposable laws, and the proof below shows it from the axioms of the scheme.

\begin{proof}[Proof of Proposition~\ref{prop:gw-ray-families}]
\emph{(1), images.} If $\widehat\mu(\omega) = \widehat\mu(c\omega)\,\widehat\rho_c(\omega)$, then at $T^{\tp}\omega$,
$\widehat{T_*\mu}(\omega) = \widehat\mu(T^{\tp}\omega) = \widehat{T_*\mu}(c\omega)\,\widehat{T_*\rho_c}(\omega)$, which is
\ref{eq:self-decomposable} for $T_*\mu$ with the residual $T_*\rho_c$.

\emph{(1), no zeros.} First an inequality: for a probability measure $\rho$ on $\RR^d$ and
$\omega \in \RR^d$,
\[
  1 - |\widehat\rho(2\omega)|^2 \le 4\bigl(1 - |\widehat\rho(\omega)|^2\bigr) .
\]
Indeed $|\widehat\rho(\omega)|^2 = \widehat\rho(\omega)\widehat\rho(-\omega)$ is the transform of
$\kappa := \rho * (-I)_*\rho$, a symmetric probability measure, so
$1 - |\widehat\rho(\omega)|^2 = \int(1 - \cos\langle\omega,x\rangle)\,\kappa(dx)$. As
$1 - \cos 2\tau = 2(1 - \cos\tau)(1 + \cos\tau) \le 4(1 - \cos\tau)$, integrating at
$\tau = \langle\omega,x\rangle$ against $\kappa$ gives the inequality. Now suppose
$\widehat\mu(\omega_1) = 0$ for some $\omega_1$, and let $t_* := \inf\{t > 0 : \widehat\mu(t\omega_1) = 0\}$. As $\widehat\mu$
is continuous with $\widehat\mu(0) = 1$, $t_* > 0$, $\omega_0 := t_*\omega_1$ has $\widehat\mu(\omega_0) = 0$, and
$\widehat\mu(t\omega_0) \ne 0$ for $0 \le t < 1$. For $c \in (0,1)$, \ref{eq:self-decomposable} at $\omega_0$ reads
$0 = \widehat\mu(c\omega_0)\,\widehat\rho_c(\omega_0)$ with $\widehat\mu(c\omega_0) \ne 0$, so $\widehat\rho_c(\omega_0) = 0$; at
$\omega_0/2$ it gives $\widehat\rho_c(\omega_0/2) = \widehat\mu(\omega_0/2)/\widehat\mu(c\omega_0/2)$, which tends to $1$ as
$c \uparrow 1$ by continuity. The inequality at $\omega = \omega_0/2$ gives
$1 \le 4(1 - |\widehat\rho_c(\omega_0/2)|^2) \to 0$, a contradiction.

\emph{(2), necessity.} By Lemma~\ref{lem:gw-kernels-from-origin}(2) the kernels from scale zero satisfy
(a)--(c). Let $\Sigma \in \PSD(d)$ and $\lambda > 1$. As $\Sigma \le_W \lambda^2\Sigma$, the cascade gives
$\mu_{\lambda^2\Sigma} = \mu_\Sigma * \rho$ with $\rho := \mu_{\Sigma,\lambda^2\Sigma}$, and (b) with
$A = \lambda I \in G$ gives $\mu_{\lambda^2\Sigma} = (\lambda I)_*\mu_\Sigma$. Applying $(\lambda^{-1}I)_*$,
$\mu_\Sigma = (\lambda^{-1}I)_*\mu_\Sigma * (\lambda^{-1}I)_*\rho$, which is \ref{eq:self-decomposable} with
$c = \lambda^{-1}$; every $c \in (0,1)$ arises. The argument reads only that $W$ contains the rays and
$G$ the dilations, so the kernels from scale zero of a cascade family are self-decomposable for every
wedge datum, not only for the ray order.

\emph{(2), sufficiency.} By (1) no $\widehat\mu_\Sigma$ has a zero, so by
Lemma~\ref{lem:gw-kernels-from-origin}(3) it suffices to give each pair $\Sigma \le_W \Sigma'$ a probability
measure $\rho$ with $\mu_{\Sigma'} = \mu_\Sigma * \rho$. For $\Sigma = 0$ take $\rho := \mu_{\Sigma'}$,
using $\mu_0 = \delta_0$. For $\Sigma' = c\,\Sigma$ with $c = 1$ take $\delta_0$; with $c > 1$, the
self-decomposability of $\mu_\Sigma$ at $c^{-1/2}$ gives $\mu_\Sigma = (c^{-1/2}I)_*\mu_\Sigma * \rho'$, and
applying $(c^{1/2}I)_*$ and (b) with $A = c^{1/2}I$, $\mu_{c\Sigma} = (c^{1/2}I)_*\mu_\Sigma = \mu_\Sigma *
(c^{1/2}I)_*\rho'$. Uniqueness and the formula for the increments are
Lemma~\ref{lem:gw-kernels-from-origin}(2)--(3).

\emph{(3).} First, the transform of a rotation-invariant law $\nu$ is radial: if $|\omega| = |\xi|$, the
reflection $Q$ in the hyperplane $(\omega - \xi)^\perp$ is orthogonal and symmetric and carries $\omega$ to $\xi$,
so $\widehat\nu(\xi) = \widehat\nu(Q^{\tp}\omega) = \widehat{Q_*\nu}(\omega) = \widehat\nu(\omega)$. With
$\vartheta(u) := \widehat\nu(\sqrt u\,e_1)$, continuous on $[0,\infty)$, this is $\widehat\nu(\omega) = \vartheta(|\omega|^2)$; and
for $\Sigma \in \PSD(d)$, $\Sigma^{1/2}$ being symmetric,
$\widehat{(\Sigma^{1/2})_*\nu}(\omega) = \widehat\nu(\Sigma^{1/2}\omega) = \vartheta(|\Sigma^{1/2}\omega|^2)
= \vartheta(\langle\omega,\Sigma\omega\rangle)$.

Now let the family be given; $(\GL(d), W)$ is a wedge datum, the ray order being invariant under every
subgroup of $\GL(d)$ (Definition~\ref{def:gw-cascade-family}). By (2), $\mu_I$ is self-decomposable, by
(a) it is not $\delta_0$, and by (b) with $A = Q \in \Orth(d)$ at $\Sigma = I$, where $Q \cdot I = I$, it
is rotation invariant. Let $\vartheta$ be its function as above. On $\PD(d)$, $\Sigma^{1/2} \in \GL(d)$ and
$\Sigma^{1/2}\cdot I = \Sigma$, so (b) with $A = \Sigma^{1/2}$ gives $\mu_\Sigma = (\Sigma^{1/2})_*\mu_I$
and $\widehat\mu_\Sigma(\omega) = \vartheta(\langle\omega,\Sigma\omega\rangle)$; this is the orbit form
$\mu_\Sigma = (\Sigma^{1/2})_*\mu_I$ of Proposition~\ref{prop:ray-families}, obtained here from (b) alone.
For $\Sigma \in \PSD(d)$, $\Sigma + \varepsilon I \in \PD(d)$ tends to $\Sigma$ as
$\varepsilon \downarrow 0$, and (c) gives $\widehat\mu_\Sigma(\omega) =
\lim_{\varepsilon\downarrow0}\vartheta(\langle\omega,(\Sigma + \varepsilon I)\omega\rangle) =
\vartheta(\langle\omega,\Sigma\omega\rangle)$, the transform of $(\Sigma^{1/2})_*\mu_I$; so
$\mu_\Sigma = (\Sigma^{1/2})_*\mu_I$ there too.

Conversely, let $\mu_I$ be rotation invariant, self-decomposable and not $\delta_0$, with function $\vartheta$,
and put $\mu_\Sigma := (\Sigma^{1/2})_*\mu_I$, so that $\widehat\mu_\Sigma(\omega) =
\vartheta(\langle\omega,\Sigma\omega\rangle)$. Then $\mu_I$ is its own member at $\Sigma = I$, as
$\vartheta(|\omega|^2) = \widehat\mu_I(\omega)$, and the conditions of (2) hold.
\begin{itemize}
\item (a), as $\vartheta(0) = 1$.
\item (b) for every $A \in \GL(d)$, as $\widehat{A_*\mu_\Sigma}(\omega) =
\vartheta(\langle A^{\tp}\omega,\Sigma A^{\tp}\omega\rangle) = \vartheta(\langle\omega,A\Sigma A^{\tp}\omega\rangle) =
\widehat\mu_{A\Sigma A^{\tp}}(\omega)$.
\item (c), as $\vartheta$ and $\Sigma \mapsto \langle\omega,\Sigma\omega\rangle$ are continuous.
\item Each $\mu_\Sigma$ is self-decomposable by (1), as an image of $\mu_I$.
\end{itemize} By (2) these are the kernels from scale zero of a
family, whose isotropic member is $\mu_I$. Two families with the isotropic member $\mu_I$ have the
kernels from scale zero $(\Sigma^{1/2})_*\mu_I$, by the first half, and so agree on the pairs by the
uniqueness in (2).
\end{proof}

\subsection{The step set of an orbit}\label{app:wedge-axis-steps}

The four proofs of this subsection decide which steps across shapes an orbit survives: the stable tail
gives necessity for the stable orbits, the ordered densities give sufficiency for every profile, the
stable order on $\PSD(d)$ carries the general linear cell, and the anisotropy budget excludes every
non-Gaussian law beyond $C_2$.

\begin{proof}[Proof of Lemma~\ref{lem:stable-tail}]
\emph{Normalisation.} Write $E(\omega)$ for the bracket in \ref{eq:stable-tail-transform} with its
sign reversed, so $\widehat\rho = e^{-E}$, and $E$ is real and even. Then $\rho(\RR^d) =
\widehat\rho(0) = 1$, and $|\widehat\rho| \le \rho(\RR^d) = 1$ gives $E \ge 0$. Put
$D := \psi_{H'} - \psi_H$, homogeneous of degree $\alpha$, so that
$E(\omega/R) = R^{-2}\langle\omega,C\omega\rangle + R^{-\alpha}D(\omega)$ for $R > 0$.

\emph{Fourier side.} Let $\phi$ be smooth with compact support in $\RR^d \setminus \{0\}$, real,
and $\widehat\phi(\omega) := \int e^{i\langle\omega,x\rangle}\phi(x)\,dx$, a Schwartz function, so
that $\phi(x) = (2\pi)^{-d}\int\widehat\phi(\omega)e^{-i\langle\omega,x\rangle}d\omega$. Put
$\phi_R(x) := \phi(x/R)$, with $\widehat{\phi_R}(\omega) = R^d\widehat\phi(R\omega)$. Integrating
the inversion formula against $\rho$ (Fubini, $\widehat\phi$ integrable and $\rho$ finite) and using
that $\widehat\rho$ is even,
\[
  \int\phi_R\,d\rho = (2\pi)^{-d}\int\widehat{\phi_R}(\omega)\,\widehat\rho(\omega)\,d\omega
  = (2\pi)^{-d}\int\widehat\phi(\zeta)\,e^{-E(\zeta/R)}\,d\zeta .
\]
Since $\int\widehat\phi = (2\pi)^d\phi(0) = 0$, the factor $e^{-E(\zeta/R)}$ may be replaced by
$e^{-E(\zeta/R)} - 1$.

\emph{The limit.} Let $R \ge 1$ and $s := E(\zeta/R) \ge 0$. Since $0 \le 1 - e^{-s} \le s$ and
$R^\alpha E(\zeta/R) = R^{\alpha-2}\langle\zeta,C\zeta\rangle + D(\zeta)$ with $R^{\alpha-2} \le 1$,
\[
  R^\alpha\bigl|e^{-E(\zeta/R)} - 1\bigr| \le R^\alpha E(\zeta/R) \le F(\zeta)
  := \|C\|\,|\zeta|^2 + |D(\zeta)| ,
\]
$\|C\|$ the operator norm. As $|D(\zeta)| \le (\|H'\|^{\alpha/2} + \|H\|^{\alpha/2})|\zeta|^\alpha \le
(\|H'\|^{\alpha/2} + \|H\|^{\alpha/2})(1 + |\zeta|^2)$ and $\widehat\phi$ is a Schwartz function,
$|\widehat\phi|F$ is integrable, and it dominates
$R^\alpha\widehat\phi(\zeta)(e^{-E(\zeta/R)} - 1)$ for $R \ge 1$. Pointwise, since $\alpha < 2$,
$R^\alpha E(\zeta/R) \to D(\zeta)$ and $E(\zeta/R) \to 0$, and $|e^{-s} - 1 + s| \le s^2$ for
$|s| \le 1$ gives
$R^\alpha\bigl|e^{-E(\zeta/R)} - 1 + E(\zeta/R)\bigr| \le R^\alpha E(\zeta/R)\cdot E(\zeta/R) \to 0$;
so $R^\alpha(e^{-E(\zeta/R)} - 1) \to -D(\zeta)$, and dominated convergence gives
\[
  \lim_{R\to\infty}R^\alpha\int\phi_R\,d\rho = -(2\pi)^{-d}\int\widehat\phi(\zeta)\,D(\zeta)\,d\zeta .
\]

\emph{Back to $x$.} Put $f := n_{H'} - n_H$, even and continuous on $\RR^d \setminus \{0\}$. By
Lemma~\ref{lem:elliptical-stable}(1), applied to $H'$ and to $H$,
$D(\zeta) = \int(1 - \cos\langle\zeta,x\rangle)f(x)\,dx$, and
$\iint|\widehat\phi(\zeta)|(1 - \cos\langle\zeta,x\rangle)(n_{H'} + n_H)(x)\,dx\,d\zeta
= \int|\widehat\phi|(\psi_{H'} + \psi_H) < \infty$ by Tonelli. So Fubini applies, and the inversion
formula at $x$ and at $-x$ gives
$\int\widehat\phi(\zeta)(1 - \cos\langle\zeta,x\rangle)\,d\zeta = (2\pi)^d\bigl[\phi(0) -
\tfrac12(\phi(x) + \phi(-x))\bigr] = -(2\pi)^d\tfrac12(\phi(x) + \phi(-x))$. Hence
\[
  -(2\pi)^{-d}\int\widehat\phi\,D = \int f(x)\,\tfrac12\bigl(\phi(x) + \phi(-x)\bigr)\,dx
  = \int f\phi\,dx ,
\]
$f$ being even. This is \ref{eq:stable-tail}.

\emph{The sign.} If $\phi \ge 0$, the left-hand side of \ref{eq:stable-tail} is a limit of
nonnegative numbers, so $\int f\phi \ge 0$. If $f(x_0) < 0$ for some $x_0 \ne 0$, then $f < 0$ on a
ball around $x_0$ not containing $0$, and a smooth $\phi \ge 0$, not identically $0$, supported in
that ball gives $\int f\phi < 0$. So $f \ge 0$.
\end{proof}

\begin{proof}[Proof of Proposition~\ref{prop:orbit-step-cone}]
\emph{Clause (1).} By C2(iii) there is an infinitely divisible $\mu_I$ with triplet
$(0, m(|x|)\,dx, 0)$. The density is even, so the odd part
$i(\sin\langle\omega,x\rangle - \langle\omega,x\rangle\mathbf 1_{\{|x| \le 1\}})$ of the integrand of
C2(i), which is integrable, integrates to $0$, and $\widehat\mu_I = e^{-\psi_I}$. The transform
determines the law, and it is radial, so $\mu_I$ is rotation invariant. The transform of $\mu_\Sigma$ at $\omega$ is
$e^{-\psi_I(\Sigma^{1/2}\omega)}$, and the substitution $x = \Sigma^{1/2}y$ gives
$\psi_I(\Sigma^{1/2}\omega) = \int(1 - \cos\langle\omega,x\rangle)\,n_\Sigma(x)\,dx$. Also
$|\Sigma^{1/2}y|^2 \wedge 1 \le \max(1, \|\Sigma\|)(|y|^2 \wedge 1)$, so
$\int(|x|^2 \wedge 1)\,n_\Sigma(x)\,dx < \infty$.

\emph{Clause (2): reduction to one direction.} Substitute $x = \Sigma^{1/2}y$:
\[
  n_\Sigma(\Sigma^{1/2}y) = (\det\Sigma)^{-1/2}m(|y|) , \qquad
  n_{\Sigma'}(\Sigma^{1/2}y) = (\det\Sigma')^{-1/2}m(\langle y,M^{-1}y\rangle^{1/2}) ,
\]
so the ordering says $m(|y|) \le (\det M)^{-1/2}m(\langle y,M^{-1}y\rangle^{1/2})$ for every $y \ne 0$. Put
$\lambda := \lambda_{\min}(M)^{1/2}$, $D := (\det M)^{1/2}$, and let $v$ be a unit eigenvector of $M$ for
$\lambda_{\min}(M)$. On the sphere $|y| = r$, $\langle y,M^{-1}y\rangle \le r^2/\lambda^2$, with equality at
$y = rv$; as $m$ is nonincreasing, the right-hand side is least there. So the ordering is
$D\,m(r) \le m(r/\lambda)$ for every $r > 0$, that is, multiplying by $r^d$,
\[
  k(r/\lambda) \ge (D/\lambda^d)\,k(r) \qquad (r > 0).
\]
\emph{The least eigenvalue is at least $1$.} $D \ge \lambda^d$, since $\det M \ge \lambda_{\min}(M)^d$.
Suppose $\lambda < 1$ and pick $r$ with $k(r) > 0$. The display gives $k(r/\lambda) \ge k(r)$ with
$r/\lambda > r$, so, $k$ being nonincreasing, $k(r/\lambda) = k(r)$, and by iteration
$k(r\lambda^{-n}) = k(r)$ for every $n \ge 1$. As $r\lambda^{-n} \to \infty$ and $k$ is nonincreasing,
$k = k(r) > 0$ on $[r,\infty)$, and then $\int_{|x| > r}m(|x|)\,dx = |S^{d-1}|\int_r^\infty k(s)\,s^{-1}ds
= \infty$, against the integrability of $m$. So $\lambda \ge 1$, and the display is
$D/\lambda^d \le R(\lambda)$, which is the stated inequality; $M \ge I$ gives
$\Sigma' - \Sigma = \Sigma^{1/2}(M - I)\Sigma^{1/2} \ge 0$.

\emph{The increment.} $f := n_{\Sigma'} - n_\Sigma \ge 0$ is even and at most $n_{\Sigma'}$, so
$f(x)\,dx$ satisfies the conditions of C2 by clause (1). C2(iii) gives an infinitely divisible $\rho$
with triplet $(0, f(x)\,dx, 0)$, and the odd part of the integrand vanishing as in clause (1),
$\widehat\rho = \exp[-\int(1 - \cos\langle\omega,x\rangle)f(x)\,dx] =
\widehat\mu_{\Sigma'}/\widehat\mu_\Sigma$, both integrals being finite.

\emph{Clause (3).} Reflexivity, transitivity and antisymmetry are those of the pointwise order of
the densities and of the Loewner order, which contains the relation. Covariance:
$|(A\Sigma A^{\tp})^{-1/2}x| = |\Sigma^{-1/2}A^{-1}x|$, so $n_{A\Sigma A^{\tp}}(x) = |\det
A|^{-1}n_\Sigma(A^{-1}x)$, and the ordering of the pair $(A\Sigma A^{\tp}, A\Sigma'A^{\tp})$ is that of
$(\Sigma,\Sigma')$ at $A^{-1}x$. A ray has $M = \lambda^2I$, $\det M = \lambda^{2d} \le
\lambda^{2d}R(\lambda)^2$. A rank-deficient step has $\lambda_{\min}(M) = 1 < \det M$, $M = I + \Sigma^{-1/2}P\Sigma^{-1/2}$ having the eigenvalue $1$
and one above $1$, while the inequality at $\lambda = 1$ reads $\det M \le
R(1)^2 = 1$. For $k = c\,r^{-\alpha}$, $k(r/\lambda)/k(r) = \lambda^\alpha$ for every $r$. If $k(0{+})$ is
finite, it is positive, $k$ being nonincreasing and not $0$. Then $R(\lambda) \le \lim_{r\downarrow0}k(r/\lambda)/k(r) = k(0{+})/k(0{+})
= 1$, while $R \ge 1$ because $k(r/\lambda) \ge k(r)$ for $\lambda \ge 1$. So $R \equiv 1$, and the inequality
becomes $\det M \le \lambda_{\min}(M)^d$. It holds only when every eigenvalue of $M$ equals $\lambda_{\min}(M)$, that
is $M = \lambda_{\min}(M)I$, a ray.
\end{proof}

\begin{proof}[Proof of Proposition~\ref{prop:similarity-stable-order-families}]
% proof of prop:similarity-stable-order-families
\emph{(1), the order.} On $\PD(d)$ the relation is $\preceq_\alpha$: a definite matrix is not $0$, and
the second case is contained in the third, every ray $\Sigma \le c\,\Sigma$, $c \ge 1$, being a step of
$\preceq_\alpha$ by Proposition~\ref{prop:stable-step-cone}(3) with $\lambda = \sqrt c$. If $\Sigma$ is definite, every
$\Sigma'$ with $\Sigma \le_\alpha \Sigma'$ is definite. If $\Sigma$ is singular and not $0$, those $\Sigma'$
are its multiples $c\Sigma$, $c \ge 1$. The properties a wedge datum requires of its order follow, whatever
$G$ is.
\begin{itemize}
\item Reflexivity is the case $c = 1$.
\item Antisymmetry: the relation is contained in the Loewner order, as $\Sigma' - 0 \ge 0$,
$c\Sigma - \Sigma = (c-1)\Sigma \ge 0$, and $\preceq_\alpha$ is contained in it
(Proposition~\ref{prop:stable-step-cone}(1)).
\item Transitivity: from $\Sigma = 0$ there is nothing to prove, from a nonzero singular $\Sigma$ the
relation runs along its multiples, and on $\PD(d)$ it is the transitivity of $\preceq_\alpha$
(Proposition~\ref{prop:stable-step-cone}(2)), the rays being steps of $\preceq_\alpha$.
\item Invariance: $A\cdot 0 = 0$, $A\cdot(c\Sigma) = c\,A\cdot\Sigma$, and $\GL(d)$ preserves $\PD(d)$ and
$\preceq_\alpha$ (Proposition~\ref{prop:stable-step-cone}(2)).
\item The rays: for $0 \le c \le c'$, $c\Sigma \le_\alpha c'\Sigma$ by the first case if $c = 0$ and by the
second, with the factor $c'/c$, otherwise.
\end{itemize}
For $\alpha \le \beta < 2$ only the third case changes: with $M := \Sigma^{-1/2}\Sigma'\Sigma^{-1/2}$,
$\Sigma \preceq_\alpha \Sigma'$ gives $\lambda_{\min}(M) \ge 1$ (Proposition~\ref{prop:stable-step-cone}(1)), so
$\lambda_{\min}(M)^{d+\beta} \ge \lambda_{\min}(M)^{d+\alpha} \ge \det M$.

\emph{(1), the chains.} From $0$ to $\Sigma'$ take $\gamma(s) := s\Sigma'$, and from $\Sigma \ne 0$ to
$c\Sigma$ take $\gamma(s) := (1 + s(c-1))\Sigma$; along both, $\gamma(s) \le_\alpha \gamma(t)$ for $s \le t$
by the rays. For $\Sigma \preceq_\alpha \Sigma'$ in $\PD(d)$ write $M = U\diag(m_1, \dots, m_d)U^{\tp}$ with
$U \in \Orth(d)$, put $M^s := U\diag(m_1^s, \dots, m_d^s)U^{\tp}$, and take the chain of
Theorem~\ref{thm:cascade-infinitely-divisible}(3)(c), $\gamma(s) := \Sigma^{1/2}M^s\Sigma^{1/2}$. It is continuous and
definite, with $\gamma(0) = \Sigma$ and $\gamma(1) = \Sigma'$. Its points are ordered by the criterion (b) of
Proposition~\ref{prop:stable-step-cone}(1). For $x = \Sigma^{1/2}y \ne 0$ put $z := U^{\tp}y$; then
\[
  \langle x, \gamma(s)^{-1}x\rangle = \sum_{i} m_i^{-s}z_i^2 =: Q_s , \qquad
  \det\gamma(s) = \det\Sigma\,(\det M)^s ,
\]
so that $n_{\gamma(s)}(x) = c_{d,\alpha}^{-1}(\det\Sigma)^{-1/2}(\det M)^{-s/2}Q_s^{-(d+\alpha)/2}$. Let
$s \le t$ and $q := \lambda_{\min}(M)$. Every $m_i$ is at least $q$, so $m_i^{-t} \le q^{-(t-s)}m_i^{-s}$ and
$Q_t \le q^{-(t-s)}Q_s$. And $\det M \le q^{d+\alpha}$ gives $(\det M)^{(t-s)/2} \le
q^{(t-s)(d+\alpha)/2}$. Together,
$(\det M)^{-s/2}Q_s^{-(d+\alpha)/2} \le (\det M)^{-t/2}Q_t^{-(d+\alpha)/2}$, which is
$n_{\gamma(s)}(x) \le n_{\gamma(t)}(x)$. So $\gamma(s) \preceq_\alpha \gamma(t)$.

\emph{(1), the thresholds.} The value $\varepsilon = 0$ is the first case, and $\Sigma \le_\alpha 0$ is
impossible for $\Sigma \ne 0$. For $\varepsilon > 0$ and $t > 0$ the two relations are steps of
$\preceq_\alpha$, with relative steps $M = \Sigma/\varepsilon$ and $M = t\Sigma^{-1}$, so that
\[
  \lambda_{\min}(\Sigma/\varepsilon) = \lambda_1/\varepsilon , \quad \det(\Sigma/\varepsilon) = \det\Sigma/\varepsilon^d ,
  \qquad
  \lambda_{\min}(t\Sigma^{-1}) = t/\lambda_d , \quad \det(t\Sigma^{-1}) = t^d/\det\Sigma .
\]
Condition (c) of Proposition~\ref{prop:stable-step-cone}(1) then reads
\[
  \lambda_1^{d+\alpha} \ge \varepsilon^\alpha\det\Sigma
  \qquad\text{and}\qquad
  t^\alpha \ge \lambda_d^{d+\alpha}/\det\Sigma ,
\]
which are the two equivalences. The ratio is $((\lambda_d/\lambda_1)^{d+\alpha})^{1/\alpha}$.

\emph{(2).} By Proposition~\ref{prop:orbit-step-cone}(1), $\mu_I$ is rotation invariant with
$\widehat\mu_I = e^{-\psi_I} > 0$, and as in the proof of Proposition~\ref{prop:gw-ray-families}(3),
$\mu_\Sigma := (\Sigma^{1/2})_*\mu_I$ has $\widehat\mu_\Sigma(\omega) = F(\langle\omega,\Sigma\omega\rangle)$
with $F(u) := \widehat\mu_I(\sqrt u\,e_1)$ continuous and positive. So the transforms have no zeros, and
(a)--(c) of Lemma~\ref{lem:gw-kernels-from-origin}(2) hold for $G = \GL(d)$: (a) as $F(0) = 1$, (b) as
$\widehat{A_*\mu_\Sigma}(\omega) = F(\langle\omega,A\Sigma A^{\tp}\omega\rangle)$, and (c) as $F$ is
continuous;
$\mu_I \ne \delta_0$ because
% CHANGED Q-0281 wave 2: blind review W2G-2
$\psi_I$ is not identically $0$, $k$ being positive near $0$. By
Lemma~\ref{lem:gw-kernels-from-origin}(3) it remains to give each pair $\Sigma \le_\alpha \Sigma'$ an
increment. From $0$: $\mu_{\Sigma'}$. Along a ray, $\Sigma' = c\Sigma$ with $c > 1$: by
Proposition~\ref{prop:orbit-step-cone}(2)--(3) the ray $I \le cI$ has an increment, $\mu_{cI} = \mu_I *
\rho$, and applying $(\Sigma^{1/2})_*$, with $\Sigma^{1/2}\sqrt c = (c\Sigma)^{1/2}$, gives $\mu_{c\Sigma} =
\mu_\Sigma * (\Sigma^{1/2})_*\rho$. On $\PD(d)$, with $M := \Sigma^{-1/2}\Sigma'\Sigma^{-1/2}$ and $q := \lambda_{\min}(M) \ge 1$, the inequality
$\det M \le q^{d+\alpha}$, condition (c) of Proposition~\ref{prop:stable-step-cone}(1), and $R(q^{1/2})^2 \ge q^\alpha$
give $\det M \le q^dR(q^{1/2})^2$, the step set of Proposition~\ref{prop:orbit-step-cone}(2), which supplies the
increment. For the stable profile $k(r) = c_{d,\beta}^{-1}r^{-\beta}$, $k(|x|)|x|^{-d}$ is $n_I(x)$ of
\ref{eq:stable-levy-density} at index $\beta$, so $\mu_I = S_\beta(I)$ by
Lemma~\ref{lem:elliptical-stable}(1)--(2), and $R(\lambda) = \lambda^\beta \ge \lambda^\alpha$ by
Proposition~\ref{prop:orbit-step-cone}(3); for the profiles of Proposition~\ref{prop:tempered-step-cone}(1)
of index $\beta$, $R(\lambda) = \lambda^\beta$ by that clause.
\end{proof}

\begin{proof}[Proof of Proposition~\ref{prop:anisotropy-budget}, clause (1)]
% proof of prop:anisotropy-budget
\emph{Infinite divisibility, and the \Levy{} measure.} Let $L_t := e^{-tB}$ and
$\mu = (L_t)_*\mu * \rho_t$. By C3, $\mu$ and every $\rho_t$ are infinitely divisible.
Let $(A_\mu, \nu, \gamma)$ be the generating triplet of $\mu$ (C2(i)). As $-I \in \Orth(d)$,
$\mu$ is symmetric, and uniqueness (C2(ii)) applied to the reflected triplet
$(A_\mu, (-I)_*\nu, -\gamma)$ gives $\gamma = 0$ and $(-I)_*\nu = \nu$, so that
$\widehat\mu(\omega) = \exp[-\frac12\langle \omega,A_\mu \omega\rangle - \int(1 - \cos\langle \omega,x\rangle)\,\nu(dx)]$.
Then $\nu \ne 0$, since otherwise $\mu$ is the Gaussian law $N(0, A_\mu)$. For $L \in \GL(d)$ this form,
read at $L^{\tp}\omega$, shows that $L_*\mu$ has the triplet $(LA_\mu L^{\tp}, L_*\nu, 0)$. The exponents of
a convolution add, so uniqueness applied to $\mu = (L_t)_*\mu * \rho_t$ gives $\nu = (L_t)_*\nu + \nu_t$
with $\nu_t$ the \Levy{} measure of $\rho_t$: hence $(L_t)_*\nu \le \nu$ for every $t > 0$. Applied to
$Q_*\mu = \mu$, $Q \in \Orth(d)$, it gives $Q_*\nu = \nu$.

\emph{Radial smoothing.} Fix a smooth $\eta$ on $\RR$ with $0 \le \eta \le 1$, $\eta = 1$ on
$[3/4, 5/4]$ and $\eta = 0$ off $(1/2, 3/2)$, and put $\nu_\eta(U) := \int\eta(s)\,\nu(s^{-1}U)\,ds$ for
Borel $U$, the average of the dilates $(sI)_*\nu$. It is invariant under $\Orth(d)$, null at the
origin and not $0$, and $(L_t)_*\nu_\eta \le \nu_\eta$, because $L_t$ commutes with $sI$. Where
$\eta(s) \ne 0$, $s < 3/2$ and $s^2|x|^2 \wedge 1 \le \frac94(|x|^2 \wedge 1)$, so
$\int(|y|^2 \wedge 1)\,\nu_\eta(dy) \le \frac94\int\eta\,\int(|x|^2 \wedge 1)\,\nu(dx) < \infty$, and
$\nu_\eta$ is finite off every ball about the origin. Substituting $r = s|x|$ and exchanging the
integrals, its image under $y \mapsto |y|$ is $K(r)\,dr$ on $(0,\infty)$, with
\[
  K(r) := \int\eta(r/|x|)\,|x|^{-1}\,\nu(dx) .
\]
For $r$ in a compact subset of $(0,\infty)$ only $x$ with $|x|$ in a compact subset of $(0,\infty)$
contribute, a set of finite $\nu$-measure, so $K$ is finite there and may be differentiated under the
integral. By polar coordinates the measure with density $h(y) := K(|y|)/(|S^{d-1}|\,|y|^{d-1})$ has the
same radial image $K(r)\,dr$, and it is invariant under $\Orth(d)$. A measure invariant under
$\Orth(d)$, null at the origin and finite off every ball about it is determined by its radial image:
on each shell $\{|x| > \varepsilon\}$ its restriction $m$ is finite and invariant, and averaging
$\widehat m(|\omega|\theta)$ over unit vectors $\theta$ against $\sigma_{d-1}$ and exchanging the integrals
gives $\widehat m(\omega) = |S^{d-1}|^{-1}\int\Phi(|\omega|,r)\,(|\cdot|_*m)(dr)$ with
$\Phi(s,r) := \int_{S^{d-1}}e^{isr\langle e_1,\theta\rangle}\,\sigma_{d-1}(d\theta)$. So $\nu_\eta$ has the
density $h$; no Haar measure on $\Orth(d)$ is needed.

\emph{The infinitesimal inequality.} Write $M(r) := K(r)/r^{d-1}$, so that $h(y) = M(|y|)/|S^{d-1}|$.
The linear change of variables $z = L_ty$, with $|\det L_t|^{-1} = e^{t\tr B}$, turns
$(L_t)_*\nu_\eta \le \nu_\eta$ into $e^{t\tr B}M(|e^{tB}z|) \le M(|z|)$ for almost every $z$, hence for
every $z \ne 0$: the set where it fails is open, both sides being continuous off the origin, and
null. Take $z = rv$ with $v$ a unit eigenvector of $\mathrm{sym}\,B$ for $\lambda$. Then
$\tfrac{d}{dt}|e^{tB}rv|$ at $t = 0$ is $r\langle v,Bv\rangle = r\lambda$, and a function of $t \ge 0$
that is largest at $t = 0$ has a nonpositive right derivative there:
\[
  \tr B\;M(r) + \lambda\,r\,M'(r) \le 0 \qquad (r > 0).
\]
Since $\det e^{-tB} = e^{-t\tr B} \to 0$, $\tr B > 0$. As $\nu_\eta \ne 0$ there is $r_0$ with
$M(r_0) > 0$.

\emph{The budget.} If $\lambda = 0$ the inequality gives $M(r_0) \le 0$, which is excluded. Otherwise
put $\kappa := \tr B/\lambda$; then $(r^\kappa M(r))' = r^{\kappa-1}(\tr B\,M(r) + \lambda rM'(r))/\lambda$,
which is $\le 0$ if $\lambda > 0$ and $\ge 0$ if $\lambda < 0$. If $\lambda < 0$, then $\kappa < 0$ and
$M(r) \ge M(r_0)r_0^\kappa r^{-\kappa} \ge M(r_0)r_0^\kappa$ for $r \ge \max(r_0, 1)$, so that
$K(r) = r^{d-1}M(r)$ is bounded below there by a positive constant and
$\int_0^\infty(r^2 \wedge 1)K(r)\,dr = \int(|y|^2 \wedge 1)\,\nu_\eta(dy)$ is infinite, which is excluded.
So $\lambda > 0$, and $M(r) \ge M(r_0)r_0^\kappa r^{-\kappa}$ for $0 < r \le r_0$, so that with
$r_1 := \min(r_0, 1)$
\[
  \int_0^{r_1}r^2K(r)\,dr \ge M(r_0)\,r_0^\kappa\int_0^{r_1}r^{d+1-\kappa}\,dr ,
\]
which is infinite if $\kappa \ge d+2$. Hence $\kappa < d + 2$, that is $\tr B < (d+2)\lambda$.
\end{proof}

\subsection{The theorem}\label{app:wedge-axis-theorem}

The theorem assembles the three subsections above along the chain of cones.
% Round-1 review (Gemini), suggestion 6 (VERIFY-round1-mathematics-paper item 13): the reason for the
% paragraph "The orders." is said here, in prose, so that the proof of record stays verbatim.
Its first paragraph derives the properties of the orders $\preceq_C$ and $\preceq^+_C$ that it needs
from the causal curves themselves. Theorem~\ref{thm:cone-orders}(4) gives them too, but
Theorem~\ref{thm:cone-orders} spends C4 through Lemma~\ref{lem:spectral-cones}(2), and the direct
derivation keeps the theorem independent of C4.

\begin{proof}[Proof of Theorem~\ref{thm:wedge-axis}]
% CHANGED Q-0281 wave 2, group H: re-transcribed from the blueprint's proof of record, the
% entry-free route (the author's decision on W2F-1): the paragraph on the orders is new, and
% (1), (2)(b), (2)(c) and (3) read the orders without C4.
\emph{The orders.} The proof reads of $\preceq_C$ and $\preceq^+_C$ only what the causal
curves give directly. By facts (c) and (a) of the proof of Proposition~\ref{prop:stable-order-generated}: for $X \in C$
the curve $t \mapsto e^{tX}$ is $C$-causal from $I$ to $e^X$, so $I \preceq_C e^X$; as $(\log c)I \in C$ for
$c \ge 1$, $\Sigma \preceq_C c\Sigma$ for $\Sigma \in \PD(d)$; and if $C \subseteq \PSD(d)$, a $C$-causal curve
$\gamma$ has $\gamma' = \gamma^{1/2}(\gamma^{-1/2}\gamma'\gamma^{-1/2})\gamma^{1/2} \ge 0$ almost everywhere, so
it is Loewner-monotone and $\preceq^+_C$ is contained in the Loewner order. A $C$-causal curve is
$C'$-causal for every cone $C'$ containing $C$, so $\preceq^+_C \subseteq \preceq^+_{C'}$. Two $C$-causal curves
$\gamma_1$, $\gamma_2$ with $\gamma_1(1) = \gamma_2(0)$ run in succession at double speed, $\gamma_1(2t)$ for
$t \le 1/2$ and $\gamma_2(2t-1)$ after, form a $C$-causal curve: its entries are absolutely continuous, each being its value at $0$ plus the
integral of the derivatives of $\gamma_1$ and $\gamma_2$ run at double speed, and its relative derivative is twice that of
$\gamma_1$ or $\gamma_2$, which lies in the cone $C$. So $\preceq_C$ is transitive, and with the rays so is
$\preceq^+_C$. For $A \in \GL(d)$, $A\gamma A^{\tp}$ is $C$-causal with $\gamma$, $C$ being $\Orth(d)$-invariant: its
relative derivative is
$Q(\gamma^{-1/2}\gamma'\gamma^{-1/2})Q^{\tp}$ with $Q := (A\gamma A^{\tp})^{-1/2}A\gamma^{1/2}$ orthogonal. So
$\preceq^+_C$ is reflexive, transitive and $\GL(d)$-invariant on $\PSD(d)$ and contains the rays.
% CHANGED Q-0281 wave 2: blind review W2H-1 (the antisymmetry, proved without C4) and W2H-3 (the
% restriction reads the invariance of the smaller order), re-transcribed from the proof of record.
Its antisymmetry is that of Theorem~\ref{thm:cone-orders}(4), and the causal curves give it directly.
(i) Every nonzero $X \in C$ has $\tr X > 0$, by the last sentence of Lemma~\ref{lem:spectral-cones}(4), as
$I \in C$. (ii) Let $\gamma$ be $C$-causal. By fact (b) of the proof of
Proposition~\ref{prop:stable-order-generated}, $\det\gamma$, a polynomial in the absolutely continuous entries of
$\gamma$, is absolutely continuous on $[0,1]$, and at almost every $t$ Jacobi's formula gives
$(\det\gamma)' = \det\gamma\,\tr Y$ with $Y := \gamma^{-1/2}\gamma'\gamma^{-1/2} \in C$, so $(\det\gamma)' \ge 0$
by (i); being the integral of its derivative, $\det\gamma$ is nondecreasing. Now let $\Sigma \preceq_C \Sigma'$
and $\Sigma' \preceq_C \Sigma$, along $C$-causal curves $\gamma_1$ from $\Sigma$ to $\Sigma'$ and $\gamma_2$ from
$\Sigma'$ to $\Sigma$. Then $\det\Sigma \le \det\Sigma' \le \det\Sigma$, so $\det\gamma_1$, nondecreasing with
equal values at $0$ and $1$, is constant, and its derivative $\det\gamma_1\,\tr Y_1$ vanishes almost
everywhere. As $\det\gamma_1 > 0$, $\tr Y_1 = 0$ almost everywhere, and so $Y_1 = 0$ almost everywhere by (i).
Hence $\gamma_1' = \gamma_1^{1/2}Y_1\gamma_1^{1/2} = 0$ almost everywhere, $\gamma_1$ is constant by fact (a),
and $\Sigma = \Sigma'$. (iii) Let $\Sigma \preceq^+_C \Sigma'$ and $\Sigma' \preceq^+_C \Sigma$. Only $0$ lies
below $0$, since $0 = c\,\Sigma''$ with $c \ge 1$ forces $\Sigma'' = 0$ and $0$ is not definite; so if one of
the two is $0$, so is the other. Otherwise both are nonzero, and both are singular or both are definite:
a matrix above a nonzero singular $\Sigma$ is a multiple $c\,\Sigma$, the third case of the definition
needing $\Sigma$ definite, and a matrix above a definite one is definite. If both are singular,
$\Sigma' = c\,\Sigma$ and $\Sigma = c'\Sigma'$ with $c, c' \ge 1$, so $cc' = 1$ and $c = 1$. If both are
definite, each of the two relations is $\preceq_C$, as $\Sigma \preceq_C c\Sigma$ for $c \ge 1$, and (ii) gives
$\Sigma = \Sigma'$. So $\preceq^+_C$ is a partial order, and $(\GL(d), \preceq^+_C)$ is a wedge datum
(Definition~\ref{def:gw-cascade-family}), $\GL(d)$ being a closed subgroup containing the dilations. A family along one
order is restricted to a smaller order by the argument of Lemma~\ref{lem:gw-kernels-from-origin}(1), which reads
of the two orders only that one contains the other and that the smaller is invariant under the smaller
group, which every order restricted to here is; and a function that agrees on the pairs
$\Sigma \preceq^+_C \Sigma'$ with a $(\GL(d), \preceq^+_C)$-cascade family is one, as (K1)--(K4) read it only
on pairs, which $\preceq^+_C$ being transitive and $\GL(d)$-invariant keeps among themselves.

\emph{(1).} Let $0 \le \alpha < \beta \le \infty$. A positive semidefinite $X$ has $\lambda_{\min}(X) \ge 0$, so
$(d+\beta)\lambda_{\min}(X) \ge (d+\alpha)\lambda_{\min}(X) \ge \tr X$ when $\beta < \infty$, and $C_\alpha
\subseteq C_\beta$ (trivially if $\beta = \infty$). Strictness: for $\beta < \infty$,
$X := \diag(1, \dots, 1, 1 + \beta)$ has $\tr X = d + \beta$ and $\lambda_{\min}(X) = 1$, so it lies in $C_\beta$
and not in $C_\alpha$; for $\beta = \infty$, $X := e_1e_1^{\tp} \in \PSD(d)$ has $\lambda_{\min} = 0 < 1 = \tr$,
so it lies in no $C_\alpha$, $\alpha < \infty$. A $C_\alpha$-causal curve is $C_\beta$-causal, so
$\preceq_{C_\alpha} \subseteq \preceq_{C_\beta}$, and $I \preceq_{C_\beta} e^X$ for the $X$ just found. But
$I \not\preceq_{C_\alpha} e^X$: by Proposition~\ref{prop:stable-order-generated}(1), $I \preceq_{C_0} e^X$ would make $e^X = cI$,
hence $X = (\log c)I \in C_0$, and for $\alpha > 0$, $I \preceq_{C_\alpha} e^X$ would give
$e^{(d+\alpha)\lambda_{\min}(X)} = \lambda_{\min}(e^X)^{d+\alpha} \ge \det e^X = e^{\tr X}$, hence $X \in C_\alpha$.
So $\preceq_{C_\alpha} \subsetneq \preceq_{C_\beta}$. For $\alpha < \infty$ and $0 \ne X \in C_\alpha$, $\tr X > 0$ and
$\lambda_{\min}(X) \ge \tr X/(d+\alpha) > 0$. The union: a nonzero $X$ in some $C_\alpha$, $\alpha < 2$, has
$\lambda_{\min}(X) > 0$ and so $(d+2)\lambda_{\min}(X) > (d+\alpha)\lambda_{\min}(X) \ge \tr X$; conversely
$(d+2)\lambda_{\min}(X) > \tr X \ge d\lambda_{\min}(X)$ gives $\lambda_{\min}(X) > 0$, so $X \ge 0$, and
$X \in C_\alpha$ with $\alpha := \tr X/\lambda_{\min}(X) - d \in [0, 2)$. The matrix $\diag(1, \dots, 1, 3)$ lies in
$C_2$, with equality, and not in the union.

\emph{(2), the kernels from scale zero.} Let a $(\GL(d), \preceq^+_C)$-cascade family be given. The ray order is
contained in $\preceq^+_C$ (its two cases are two of the three), so the family restricts to a
$(\GL(d), \text{rays})$-family with the same kernels from scale zero, and
Proposition~\ref{prop:gw-ray-families}(3) gives $\mu_\Sigma = (\Sigma^{1/2})_*\mu_I$ with $\mu_I$ rotation invariant,
self-decomposable and not $\delta_0$, and $\widehat\mu_\Sigma(\omega) = F(\langle\omega,\Sigma\omega\rangle)$ with
$F(u) = \widehat\mu_I(\sqrt u\,e_1)$ continuous, $F(0) = 1$. $F$ has no zero (Proposition~\ref{prop:gw-ray-families}(1)) and is
real ($\mu_I$ is invariant under $-I$), so $F > 0$, and $g := -\log F$ is continuous, $g(0) = 0$, $g \ge 0$ (as
$|F| \le 1$), and $g \not\equiv 0$ (as $\mu_I \ne \delta_0$). For $X \in C$, $I \preceq^+_C e^X$, and (K1) at
$0 \preceq^+_C I \preceq^+_C e^X$ gives $\mu_{e^X} = \mu_{I,e^X} * \mu_I$, that is
\begin{equation}\label{eq:wedge-axis-step}
  e^{-g(\langle\omega, e^X\omega\rangle)} = e^{-g(|\omega|^2)}\,\widehat\mu_{I,e^X}(\omega)
  \qquad (\omega \in \RR^d).
\end{equation}

\emph{(2)(a).} Let $X \in C$ have an eigenvalue $\lambda < 0$ with unit eigenvector $v$, and $m := e^\lambda < 1$.
At $\omega = uv$, \ref{eq:wedge-axis-step} and $|\widehat\mu_{I,e^X}| \le 1$ give $g(mu^2) \ge g(u^2)$, so $g(x) \le g(mx) \le
\dots \le g(m^nx) \to g(0) = 0$ for every $x \ge 0$, against $g \not\equiv 0$.

\emph{(2)(b).} Let $0 \ne X \in C$ be singular and $P := e^X - I$, positive semidefinite with $\rank P =
\rank X$, since its eigenvalues $e^{x_i} - 1$ vanish exactly where those of $X$ do; so $0 < \rank P < d$. \ref{eq:wedge-axis-step}
is \ref{eq:one-step} with $\Sigma = I$ and $\rho = \mu_{I,I+P}$, and Theorem~\ref{thm:rank-deficient-step} gives
$g(u) = bu$ with $b = g(1)$, $b > 0$ as $g \not\equiv 0$. So $\widehat\mu_\Sigma(\omega) =
e^{-b\langle\omega,\Sigma\omega\rangle}$ for every $\Sigma$. For $\Sigma \preceq^+_C \Sigma'$, (K1) at
$0 \preceq^+_C \Sigma \preceq^+_C \Sigma'$ gives $\widehat\mu_{\Sigma'} = \widehat\mu_{\Sigma,\Sigma'}\,\widehat\mu_\Sigma$,
and $\widehat\mu_\Sigma$ has no zeros, so $\widehat\mu_{\Sigma,\Sigma'} = e^{-b\langle\omega,(\Sigma'-\Sigma)\omega\rangle}$,
the transform of the affine Gaussian kernel ($\Sigma \le \Sigma'$ in the Loewner order, as
$C \subseteq \PSD(d)$): the family is the restriction of the affine Gaussian scale space. Conversely, the affine
Gaussian scale space with constant $b > 0$ is an affine cascade family (Theorem~\ref{thm:affine-rigidity}), that is a
$(\GL(d), \text{Loewner})$-family; $\preceq^+_C$ is contained in the Loewner order as $C \subseteq \PSD(d)$, so
its restriction is a $(\GL(d), \preceq^+_C)$-family, and so is every function that agrees with it on the pairs.

\emph{(2)(c), the number $\beta(C)$.} Now $C \subseteq \PSD(d)$ has no nonzero singular member, so its nonzero
members are positive definite. On them $f(X) := \tr X/\lambda_{\min}(X) - d$ is continuous,
nonnegative and invariant under positive scaling. The slice $\{X \in C : \tr X = 1\}$ is closed, and bounded,
since a positive semidefinite $X$ with $\tr X = 1$ has entries in $[-1, 1]$; it contains $I/d$, and every
$X \in C \setminus \{0\}$ has the multiple $X/\tr X$ in it. So $f$ attains on the slice a largest value
$\beta(C) \ge 0$, which is its largest value on $C \setminus \{0\}$. For $\alpha \in [0, \infty)$,
$C \subseteq C_\alpha$ means $(d+\alpha)\lambda_{\min}(X) \ge \tr X$, that is $\alpha \ge f(X)$, for every nonzero
$X \in C$, which is $\alpha \ge \beta(C)$.

\emph{(2)(c), the stable orbits.} If $\gamma \ge \beta(C)$, then $C \subseteq C_\gamma$, so every $C$-causal
curve is $C_\gamma$-causal and $\preceq^+_C \subseteq \preceq^+_{C_\gamma}$; by Proposition~\ref{prop:stable-order-generated}(1),
$\preceq^+_{C_\gamma}$ is the order $\le_\gamma$ of Proposition~\ref{prop:similarity-stable-order-families}(1), and that
proposition's clause (2), read with its $\alpha$ equal to $\gamma$ and the stable law of index $\gamma$, gives a
$(\GL(d), \le_\gamma)$-family with kernels $(\Sigma^{1/2})_*S_\gamma(I)$, which restricts. Conversely, if these
are the kernels of a $(\GL(d), \preceq^+_C)$-family and $X \in C$, then $\mu_{e^X} = (e^{X/2})_*S_\gamma(I) =
S_\gamma(e^X)$ (Lemma~\ref{lem:elliptical-stable}(2)) is $S_\gamma(I) * \mu_{I,e^X}$, and
Proposition~\ref{prop:stable-step-cone}(1), (a)$\Rightarrow$(c), with $\Sigma = I$ and $\Sigma' = e^X$, gives
$\lambda_{\min}(e^X)^{d+\gamma} \ge \det e^X$, that is $X \in C_\gamma$ (Proposition~\ref{prop:stable-order-generated}(1)).
So $C \subseteq C_\gamma$ and $\gamma \ge \beta(C)$.

\emph{(2)(c), the non-Gaussian families.} If $\beta(C) < 2$, take any index $\gamma$ with
$\beta(C) \le \gamma < 2$ and $\gamma > 0$, for instance $\gamma := \max(\beta(C), 1)$. The
stable orbit of index $\gamma$ is a family, and $S_\gamma(I)$ is not Gaussian, its transform
$e^{-|\omega|^\gamma}$ having modulus $e^{-t^\gamma}$ at $\omega = te_1$, which is $e^{-at^2}$ for no $a \ge 0$.
Conversely let a family have a non-Gaussian $\mu_I$, and let $0 \ne X \in C$ and $t \ge 0$. By (K1), as in \ref{eq:wedge-axis-step} with
$tX \in C$ in place of $X$, and the form of the kernels, $(e^{tX/2})_*\mu_I = \mu_{e^{tX}} = \mu_I * \mu_{I,e^{tX}}$; applying
$(e^{-tX/2})_*$, $\mu_I = (e^{-tB})_*\mu_I * (e^{-tB})_*\mu_{I,e^{tX}}$ with $B := X/2$, so $\mu_I$ is
$B$-self-decomposable in the sense of Proposition~\ref{prop:stable-exponent-cone}. $B$ is positive definite, so
$e^{-tB} \to 0$, and $\mu_I$ is rotation invariant and not Gaussian; Proposition~\ref{prop:anisotropy-budget}(1) gives
$\tr B < (d+2)\lambda_{\min}(B)$, that is $f(X) < 2$. As the largest value of $f$ is attained, $\beta(C) < 2$.
If $\beta(C) \ge 2$, then every family has a Gaussian $\mu_I$. Its image under $x \mapsto \langle e_1, x\rangle$
is a Gaussian law on $\RR$, of some variance $v \ge 0$, so $e^{-g(t^2)} = |\widehat\mu_I(te_1)| = e^{-vt^2/2}$ and
$g(u) = bu$ with $b := v/2$, and $b > 0$ as $g \not\equiv 0$; the end of (2)(b) shows that the
family is the restricted affine Gaussian scale space; these exist, by the converse there, which used only
$C \subseteq \PSD(d)$.

\emph{(3).} By Proposition~\ref{prop:stable-order-generated}(1), $\Sigma \preceq_{C_0} \Sigma'$ on $\PD(d)$ means $\Sigma' = c\Sigma$
with $c \ge 1$, so by the definition of $\preceq^+$ the order $\preceq^+_{C_0}$ is the ray order, and
Proposition~\ref{prop:gw-ray-families}(3) is the description of $\mathfrak F_0$. For $\alpha \le \beta$,
$\preceq^+_{C_\alpha} \subseteq \preceq^+_{C_\beta}$ by (1), and a family restricts; so
$\mathfrak F_\beta \subseteq \mathfrak F_\alpha$. For $0 \le \alpha < \infty$, $C_\alpha$ satisfies the
hypotheses of (2)(c), as its nonzero members are definite by (1), and $\beta(C_\alpha) = \alpha$, since
$f \le \alpha$ on $C_\alpha$ with equality at $\diag(1, \dots, 1, 1+\alpha)$. The clauses of (3) follow,
range by range.
\begin{itemize}
\item $\alpha = 0$: $\mathfrak F_0$ is described above.
\item $0 < \alpha < 2$: the stable clause is (2)(c), and the profile clause is
Proposition~\ref{prop:similarity-stable-order-families}(2), its order $\le_\alpha$ being $\preceq^+_{C_\alpha}$.
\item $2 \le \alpha < \infty$: (2)(c) with $\beta(C_\alpha) \ge 2$ gives the Gaussians.
\item $\alpha = \infty$: $C_\infty = \PSD(d)$ contains $e_1e_1^{\tp}$, so (2)(b) gives the Gaussians.
\end{itemize}
Strictness for $\alpha < \beta \le 2$: let $\gamma := \alpha$ if $\alpha > 0$ and $\gamma := \beta/2$ if
$\alpha = 0$, so $0 < \gamma < \beta$ and $\gamma \ge \alpha$. By (2)(c) with $\beta(C_\alpha) = \alpha$,
$S_\gamma(I) \in \mathfrak F_\alpha$; and $S_\gamma(I) \notin \mathfrak F_\beta$, by (2)(c) with $\beta(C_\beta) =
\beta > \gamma$ if $\beta < 2$, and because $\mathfrak F_2$ contains only Gaussian laws if $\beta = 2$.
\end{proof}
